\documentclass[11pt,a4paper]{article}
\usepackage[T1]{fontenc}
\usepackage[utf8]{inputenc}
\usepackage{lmodern}
\usepackage{microtype}
\usepackage{amsmath,amssymb,amsthm}
\usepackage{booktabs}
\usepackage{array}
\usepackage{longtable}
\usepackage{enumitem}
\usepackage[margin=2.6cm]{geometry}
\usepackage{xcolor}
\usepackage[colorlinks=true,linkcolor=blue!60!black,citecolor=blue!60!black,urlcolor=blue!60!black]{hyperref}

\newtheorem{theorem}{Theorem}[section]
\newtheorem{proposition}[theorem]{Proposition}
\newtheorem{lemma}[theorem]{Lemma}
\newtheorem{corollary}[theorem]{Corollary}
\newtheorem{conjecture}[theorem]{Conjecture}
\newtheorem{question}[theorem]{Question}
\theoremstyle{definition}
\newtheorem{definition}[theorem]{Definition}
\newtheorem{example}[theorem]{Example}
\theoremstyle{remark}
\newtheorem{remark}[theorem]{Remark}
\newtheorem{observation}[theorem]{Observation}

\newcommand{\Cay}{\operatorname{Cay}}
\newcommand{\Aut}{\operatorname{Aut}}
\newcommand{\ord}{\operatorname{ord}}
\newcommand{\PSL}{\operatorname{PSL}}
\newcommand{\PGL}{\operatorname{PGL}}
\newcommand{\PSigmaL}{\operatorname{P\Sigma L}}
\newcommand{\GF}{\operatorname{GF}}
\newcommand{\ZZ}{\mathbb{Z}}
\newcommand{\X}[3]{X(#1;#2,#3)}
\newcommand{\Tr}{\operatorname{T}}

\title{The Alspach--Zhang conjecture:\\ every cubic Cayley graph is $3$-edge-colourable\\[1ex]\large A survey}
\author{}
\date{September 2026}

\begin{document}
\maketitle

\begin{abstract}
In 1992 Brian Alspach and Cun-Quan Zhang conjectured that every Cayley graph of valency at least two admits a nowhere-zero $4$-flow.  By theorems of Mader and Jaeger this is equivalent to the statement that every connected cubic Cayley graph is $3$-edge-colourable, or, in the language introduced by Nedela and \v{S}koviera, that there are no \emph{Cayley snarks}.  The conjecture is implied by the folklore conjecture that every Cayley graph is hamiltonian, and it is known to hold for solvable groups (Alspach, Liu and Zhang), for all cubic graphs admitting a solvable vertex-transitive group of automorphisms other than the Petersen graph (Poto\v{c}nik), and for all cubic Cayley graphs on at most $1280$ vertices (Poto\v{c}nik, Spiga and Verret).  A minimal counterexample, if one exists, is a Cayley graph on a non-abelian simple group or on a group with a unique non-trivial proper normal subgroup, which has index two (Nedela and \v{S}koviera).  We survey these results, give short self-contained proofs of the elementary reductions, record several reformulations (in terms of arc-regular cubic graphs, of bi-Cayley graphs, and of small girth), prove a quotient lemma for arbitrary subgroups (with semi-edges), describe the $3$-edge-colourings of $\Cay(G,\{a,x,x^{-1}\})$ as $a$-invariant functions $G\to\mathbb{Z}_3$ with a local admissibility condition along the $x$-cycles, and prove that the index-two case is always colourable when $\ord(x)=3$.  We also analyse the new $\mathbb F_2^3$ construction in the 2026 proof of the cycle double cover theorem, extracting an exact eight-vertex palette reformulation, constructing generator-separated rank-three flows, and reducing the potentially non-circular part of the method to twelve explicit augmented affine systems.  Finally, we report a computer verification of the conjecture for all cubic Cayley graphs on 85 non-solvable groups of order up to $300\,696$, including every group permitted by the Nedela--\v{S}koviera reduction of order less than $352\,440$; consequently a smallest Cayley snark, if one exists, has more than $352\,439$ vertices.
\end{abstract}

\tableofcontents

\section{Introduction}

A \emph{Cayley graph} $\Cay(G,S)$ of a finite group $G$ with respect to a subset $S\subseteq G\setminus\{1\}$ with $S=S^{-1}$ has vertex set $G$, with $g$ adjacent to $gs$ for every $s\in S$.  Cayley graphs are the most symmetric graphs one can build from a group: $G$ itself acts regularly on the vertices by left multiplication.  Many of the classical ``hard'' cubic graphs of graph theory---the Petersen graph, the Coxeter graph, the snarks---are highly symmetric, yet none of them is a Cayley graph.  The conjecture surveyed here asserts that this is not an accident.

\begin{conjecture}[Alspach--Zhang, 1992]\label{conj:AZ}
Every connected cubic Cayley graph is $3$-edge-colourable.
\end{conjecture}

The conjecture was formulated by Alspach and Zhang at the Louisville workshop on Hamilton cycles in 1992, in the language of flows: \emph{every Cayley graph of valency at least two admits a nowhere-zero $4$-flow} (see \cite[\S1]{LiZhou16}, \cite[\S1]{ZhangZhou24}).  It first appeared in print in the paper of Alspach, Liu and Zhang \cite{ALZ96}, who proved it for solvable groups.  Since a connected vertex-transitive graph of valency $k$ is $k$-edge-connected (Mader \cite{Mader71}) and every $4$-edge-connected graph has a nowhere-zero $4$-flow (Jaeger \cite{Jaeger79}), the flow version reduces at once to the cubic case, where a nowhere-zero $4$-flow is the same thing as a $3$-edge-colouring (Tutte \cite{Tutte54}); see Proposition~\ref{prop:flow} below.  Nedela and \v{S}koviera \cite{NS01} coined the term \emph{Cayley snark} for a cubic Cayley graph that is not $3$-edge-colourable; Conjecture~\ref{conj:AZ} says that Cayley snarks do not exist.

The conjecture sits between two much older problems.

\begin{itemize}[leftmargin=2em]
\item \emph{Snarks and the Four Colour Theorem.}  Tait observed in 1880 that the Four Colour Theorem is equivalent to the statement that every bridgeless planar cubic graph is $3$-edge-colourable.  Non-planar bridgeless cubic graphs need not be: the Petersen graph is the smallest example, and the non-trivial examples are called snarks.  A large part of the theory of snarks consists in showing that highly structured cubic graphs are not snarks; the Alspach--Zhang conjecture is the statement of this kind for the strongest possible symmetry assumption compatible with the Petersen graph not being an exception.
\item \emph{Hamilton cycles in vertex-transitive graphs.}  Lov\'asz asked in 1969 whether every connected vertex-transitive graph has a Hamilton path, and a folklore conjecture asserts that every connected Cayley graph on at least three vertices has a Hamilton cycle \cite{CurranGallian96,KutnarMarusic09}.  A cubic Cayley graph has even order (Lemma~\ref{lem:structure}), so a Hamilton cycle in it is an even cycle whose edges can be coloured alternately, the remaining perfect matching receiving the third colour.  Thus \emph{Conjecture~\ref{conj:AZ} is a consequence of the hamiltonicity conjecture for cubic Cayley graphs}, and every hamiltonicity result for cubic Cayley graphs is a result about Conjecture~\ref{conj:AZ}.  Only four connected vertex-transitive graphs without a Hamilton cycle are known---the Petersen graph, the Coxeter graph, and their truncations---and none of them is a Cayley graph.
\end{itemize}

\subsection*{What is known}
Table~\ref{tab:status} summarises the state of the conjecture.  The three main theorems are those of Alspach, Liu and Zhang (solvable groups), Nedela and \v{S}koviera (structure of a minimal counterexample) and Poto\v{c}nik (solvable vertex-transitive groups); the strongest computational evidence comes from the census of cubic vertex-transitive graphs of Poto\v{c}nik, Spiga and Verret.  We add to this a computation of our own (Section~\ref{sec:computation}) which pushes the verification well beyond the range of the census for the groups singled out by the Nedela--\v{S}koviera reduction: a smallest Cayley snark would have more than $352\,439$ vertices.

\begin{table}[ht]
\centering
\small
\begin{tabular}{@{}p{0.55\textwidth}p{0.38\textwidth}@{}}
\toprule
Class of cubic Cayley graphs $\Cay(G,S)$ & Status / reference\\
\midrule
$S$ consists of three involutions & trivial (colour by generator), \S\ref{sec:easy}\\
$S=\{a,x,x^{-1}\}$ with $\ord(x)$ even & trivial, \S\ref{sec:easy}\\
$G$ abelian, dihedral, nilpotent, supersolvable, of odd-order-times-two, \dots & all solvable: \cite{ALZ96}\\
$G$ solvable & Alspach--Liu--Zhang \cite{ALZ96}\\
any cubic graph with a solvable vertex-transitive group of automorphisms, other than the Petersen graph & Poto\v{c}nik \cite{Potocnik04}\\
$|G|\le 1280$ & census \cite{PSV13} (all are hamiltonian)\\
$G=\langle a,x\rangle$ with $a^2=x^s=(ax)^3=1$ (``$(2,s,3)$-Cayley graphs'') & hamiltonian for $s$ odd \cite{GM07,GKMM12}; trivial for $s$ even\\
$G$ has a normal subgroup containing $a$, or a normal subgroup avoiding $a$ and $x$ whose quotient graph is colourable & Lemmas~\ref{lem:cover}, \ref{lem:a-in-N}\\
minimal counterexample & $G$ non-abelian simple, or $G$ has a unique non-trivial proper normal subgroup $T$, $|G:T|=2$, $T$ simple or $T\cong S\times S$ \cite{NS01}\\
girth $4$; girth $3$ or $5$ unless $\ord(x)=3$ or $5$ & Proposition~\ref{prop:girth}\\
$G$ has a subgroup $H$ avoiding the conjugacy class of $x$ such that the pregraph $H\backslash X$ is $3$-edge-colourable & Lemma~\ref{lem:pregraph}\\
$\ord(x)=3$ and $G$ has a normal subgroup of index two & Proposition~\ref{prop:index2-three}\\
$\ord(x)=3$ and $|G|\le 30\,000$ & Conder's census \cite{Conder11}, Corollary~\ref{cor:conder}\\
every group admissible for a smallest Cayley snark of order $<352\,440$, and further non-solvable groups of order $\le 300\,696$ & Section~\ref{sec:computation}\\
\bottomrule
\end{tabular}
\caption{Known cases of Conjecture~\ref{conj:AZ}.  Throughout, $a$ denotes an involution and $x$ an element of odd order $\ge 3$ generating $G$ together with $a$; by \S\ref{sec:easy} this is the only case that matters.}
\label{tab:status}
\end{table}

\subsection*{Organisation}
Section~\ref{sec:prelim} fixes terminology and recalls the classical facts about $3$-edge-colourings, flows, snarks and coverings that are used throughout.  Section~\ref{sec:easy} disposes of the trivial cases and reduces the conjecture to Cayley graphs $\Cay(G,\{a,x,x^{-1}\})$ with $a$ an involution and $x$ of odd order; it also contains several elementary reformulations that we have not seen stated explicitly in the literature (arc-regular cubic graphs, bi-Cayley graphs, small girth), and settles the order-three case when $G$ has a normal subgroup of index two.  Section~\ref{sec:new} extends the quotient lemma to arbitrary subgroups (the quotient is then a pregraph with semi-edges), gives a cyclewise colouring criterion for such quotients, describes the $3$-edge-colourings of $\Cay(G,\{a,x,x^{-1}\})$ as $a$-invariant functions $G\to\mathbb{Z}_3$ satisfying a local condition along the $x$-cycles, and derives a consecutive-$2$-factor criterion for the index-two case.  Section~\ref{sec:solvable} treats solvable groups and Poto\v{c}nik's generalisation to solvable vertex-transitive groups.  Section~\ref{sec:minimal} presents the Nedela--\v{S}koviera reduction with a self-contained proof.  Section~\ref{sec:hamilton} surveys the hamiltonicity route and the Cayley-map approach of Glover, Kutnar, Malni\v{c}, Maru\v{s}i\v{c} and Hujdurovi\'c.  Section~\ref{sec:computation} describes the computational evidence, including our own.  Section~\ref{sec:related} discusses related flow problems and reduces the potentially useful part of the 2026 cycle-double-cover method to augmented affine systems; Section~\ref{sec:open} lists open questions.

\section{Preliminaries}\label{sec:prelim}

\subsection{Cayley graphs}
Let $G$ be a finite group and $S\subseteq G\setminus\{1\}$ a subset closed under inversion.  The Cayley graph $\Cay(G,S)$ has vertex set $G$ and edge set $\{\{g,gs\}: g\in G,\ s\in S\}$.  It is a simple $|S|$-regular graph; it is connected if and only if $S$ generates $G$.  Every edge $\{g,gs\}$ carries the label $\{s,s^{-1}\}$, and left multiplication by any $h\in G$ is a label-preserving automorphism, so $G$ acts regularly on the vertices.  Conversely, a graph admitting a group of automorphisms acting regularly on its vertices is a Cayley graph on that group.  All graphs in this survey are finite; ``Cayley graph'' will always mean a connected one.

\subsection{Cubic Cayley graphs}
Let $\Cay(G,S)$ be cubic, so $|S|=3$.  Since $S=S^{-1}$ and $|S|$ is odd, $S$ contains an involution.  Either all three elements of $S$ are involutions, or $S=\{a,x,x^{-1}\}$ with $a^2=1$ and $\ord(x)\ge 3$.  We use the notation
\[
\X{G}{a}{x} := \Cay\bigl(G,\{a,x,x^{-1}\}\bigr),\qquad a^2=1,\ \ \ord(x)\ge 3,\ \ G=\langle a,x\rangle .
\]
In $\X{G}{a}{x}$ the $a$-edges form a perfect matching $M_a$, and the $x$-edges form a $2$-factor $F_x$ consisting of $|G|/\ord(x)$ cycles of length $\ord(x)$ (the left cosets of $\langle x\rangle$).

\begin{lemma}\label{lem:structure}
A connected cubic Cayley graph has even order and is either $\Cay(G,\{a,b,c\})$ with three involutions $a,b,c$, or $\X{G}{a}{x}$ as above.
\end{lemma}
\begin{proof}
The generating set contains an involution, so $|G|$ is even; the dichotomy was explained above.
\end{proof}

\subsection{Tait colourings, even 2-factors and flows}
A \emph{$3$-edge-colouring} (or Tait colouring) of a cubic graph $X$ is a colouring of the edges with three colours such that the three edges at each vertex receive distinct colours; equivalently, a partition of $E(X)$ into three perfect matchings.  The following reformulation is used constantly.

\begin{lemma}\label{lem:even2factor}
A cubic graph is $3$-edge-colourable if and only if it has a $2$-factor all of whose cycles are even.
\end{lemma}
\begin{proof}
Two colour classes of a $3$-edge-colouring form a $2$-factor of even cycles; conversely, colour the cycles of an even $2$-factor alternately and give the complementary perfect matching the third colour.
\end{proof}

In particular a cubic graph with a Hamilton cycle is $3$-edge-colourable if and only if its order is even, and a bipartite cubic graph is $3$-edge-colourable (K\H{o}nig).

A \emph{nowhere-zero $k$-flow} on a graph is an orientation together with an assignment of integers from $\{\pm1,\dots,\pm(k-1)\}$ to the edges satisfying Kirchhoff's law at every vertex.  Tutte \cite{Tutte54,Tutte66} showed that a cubic graph admits a nowhere-zero $4$-flow if and only if it is $3$-edge-colourable, and that a planar graph admits a nowhere-zero $4$-flow if and only if its dual is $4$-vertex-colourable; his $5$-flow conjecture (every bridgeless graph has a nowhere-zero $5$-flow) and $3$-flow conjecture (every $4$-edge-connected graph has a nowhere-zero $3$-flow) remain open.  Two theorems are needed to relate the flow form of the Alspach--Zhang conjecture to its colouring form.

\begin{theorem}[Jaeger \cite{Jaeger79}]\label{thm:jaeger}
Every $4$-edge-connected graph admits a nowhere-zero $4$-flow.
\end{theorem}

\begin{theorem}[Mader \cite{Mader71}]\label{thm:mader}
A connected vertex-transitive graph of valency $k$ is $k$-edge-connected.
\end{theorem}

\begin{proposition}\label{prop:flow}
The following statements are equivalent.
\begin{enumerate}[label=\textup{(\roman*)}]
\item Every connected Cayley graph of valency at least $2$ admits a nowhere-zero $4$-flow.
\item Every connected cubic Cayley graph admits a nowhere-zero $4$-flow.
\item Every connected cubic Cayley graph is $3$-edge-colourable \textup{(Conjecture~\ref{conj:AZ})}.
\end{enumerate}
\end{proposition}
\begin{proof}
A connected $2$-regular graph is a cycle and has a nowhere-zero $2$-flow.  A Cayley graph of valency $k\ge4$ is $k$-edge-connected by Theorem~\ref{thm:mader} and hence has a nowhere-zero $4$-flow by Theorem~\ref{thm:jaeger}.  So (ii)$\Rightarrow$(i), the converse is trivial, and (ii)$\Leftrightarrow$(iii) is Tutte's theorem.
\end{proof}

We also recall that a $3$-edge-colourable cubic graph has a cycle double cover (take the three $2$-factors formed by pairs of colour classes).  Thus Conjecture~\ref{conj:AZ} implied the cycle double cover conjecture for cubic Cayley graphs; the latter was proved in full generality in 2026 \cite{OpenAICDC26,Geelen26,Oum26}.  We analyse whether that proof can conversely help with Conjecture~\ref{conj:AZ} in \S\ref{sec:cdc-palette}.

\subsection{Snarks and cyclic connectivity}
A \emph{snark} is usually defined as a cyclically $4$-edge-connected cubic graph of girth at least $5$ which is not $3$-edge-colourable; the connectivity and girth conditions exclude examples obtained trivially from smaller ones.  Following \cite{NS01} we drop these conditions and call \emph{any} cubic Cayley graph that is not $3$-edge-colourable a \emph{Cayley snark}.  For vertex-transitive cubic graphs the distinction is small:

\begin{theorem}[Nedela--\v{S}koviera \cite{NS95}]\label{thm:cyclic}
The cyclic edge-connectivity of a connected cubic vertex-transitive graph equals its girth.
\end{theorem}

Thus a Cayley snark of girth at least $5$ would be a snark in the classical sense, and Jaeger's conjecture \cite{Jaeger85,JaegerSwart80} that there are no cyclically $7$-edge-connected snarks (still open; the analogous conjecture for girth was refuted by Kochol \cite{Kochol96}) would imply Conjecture~\ref{conj:AZ} for all cubic Cayley graphs of girth at least $7$.  We return to this in \S\ref{sec:girth}.

\subsection{Coverings and quotients}\label{sec:cover}
A surjective graph homomorphism $\pi\colon \tilde X\to X$ between (multi)graphs is a \emph{covering projection} if it maps the set of edges incident with each vertex $\tilde v$ bijectively onto the set of edges incident with $\pi(\tilde v)$.  If a group $N$ acts semiregularly on $\tilde X$ (every non-identity element acts without fixed vertices and without inverting edges) and the quotient $\tilde X/N$ has no loops or semi-edges, then $\tilde X\to\tilde X/N$ is a covering projection, called a \emph{regular $N$-covering}.

\begin{lemma}\label{lem:lift}
If $\pi\colon\tilde X\to X$ is a covering projection of cubic multigraphs and $X$ is $3$-edge-colourable, then so is $\tilde X$.
\end{lemma}
\begin{proof}
Colour each edge $\tilde e$ of $\tilde X$ with the colour of $\pi(\tilde e)$.  At each vertex $\tilde v$ the three edges map bijectively onto the three (distinctly coloured) edges at $\pi(\tilde v)$.
\end{proof}

For Cayley graphs the relevant quotients come from normal subgroups.  Let $N\trianglelefteq G$ and write $\bar g=gN$.

\begin{lemma}[quotient lemma]\label{lem:cover}
Let $X=\X{G}{a}{x}$ with $\ord(x)$ odd and let $1\ne N\trianglelefteq G$ be a proper normal subgroup with $a\notin N$ and $x\notin N$.  Then $\bar G=G/N$ is generated by the involution $\bar a$ and the element $\bar x$ of odd order at least $3$, $\bar X=\X{\bar G}{\bar a}{\bar x}$ is a connected cubic Cayley graph on a group of order $|G|/|N|$, and $g\mapsto\bar g$ is a regular $N$-covering $X\to\bar X$.  Consequently, if $\bar X$ is $3$-edge-colourable then so is $X$.
\end{lemma}
\begin{proof}
$\bar a\ne1$ since $a\notin N$, and $\ord(\bar x)$ divides the odd number $\ord(x)$ and is $>1$ since $x\notin N$; hence $\ord(\bar x)\ge3$, $\bar x\ne\bar x^{-1}$ and $\bar x\ne\bar a$.  So $\{\bar a,\bar x,\bar x^{-1}\}$ is an inverse-closed $3$-subset of $\bar G\setminus\{1\}$ generating $\bar G$.  The map $g\mapsto \bar g$ sends the three edges $\{g,ga\},\{g,gx\},\{g,gx^{-1}\}$ at $g$ to the three distinct edges $\{\bar g,\bar g\bar a\},\{\bar g,\bar g\bar x\},\{\bar g,\bar g\bar x^{-1}\}$ at $\bar g$, and $N$ acts semiregularly by left multiplication with quotient $\bar X$.  Apply Lemma~\ref{lem:lift}.
\end{proof}

When $N$ contains $a$ or $x$ the quotient degenerates (the $a$-edges become semi-edges, or the $x$-edges become loops).  These two cases behave very differently: the first is harmless (Lemma~\ref{lem:a-in-N}), the second is the heart of the problem (\S\ref{sec:bicayley}).

\section{The easy cases and some reformulations}\label{sec:easy}

\subsection{Trivial colourings}
\begin{proposition}\label{prop:trivial}
Let $X$ be a connected cubic Cayley graph.
\begin{enumerate}[label=\textup{(\alph*)}]
\item If $X=\Cay(G,\{a,b,c\})$ with $a,b,c$ involutions, then colouring each edge by its generator is a $3$-edge-colouring.
\item If $X=\X{G}{a}{x}$ with $\ord(x)$ even, then $F_x$ is an even $2$-factor and $X$ is $3$-edge-colourable.
\item If $G$ is abelian, then $X$ is $3$-edge-colourable.
\end{enumerate}
\end{proposition}
\begin{proof}
(a) and (b) are immediate from Lemma~\ref{lem:even2factor}.  For (c), by (a) and (b) we may assume $X=\X{G}{a}{x}$ with $\ord(x)$ odd; as $G$ is abelian, $G=\langle x\rangle\times\langle a\rangle$ and $X$ is the prism $C_{\ord(x)}\square K_2$, which has a Hamilton cycle (use two $a$-edges and all but one edge of each of the two $x$-cycles).
\end{proof}

\begin{corollary}\label{cor:reduction}
Conjecture~\ref{conj:AZ} is equivalent to the following statement: for every finite group $G=\langle a,x\rangle$ generated by an involution $a$ and an element $x$ of odd order $s\ge3$, the graph $\X{G}{a}{x}$ is $3$-edge-colourable.
\end{corollary}

From now on we therefore assume that $X=\X{G}{a}{x}$ with $s=\ord(x)$ odd, $s\ge 3$.  Note that $X$ then contains the odd cycles of $F_x$, so $X$ is never bipartite, and a $3$-edge-colouring must use a $2$-factor different from the ``obvious'' one.  In terms of the group, $X$ is $3$-edge-colourable if and only if there is an $a$-invariant subset $M\subseteq G$ and a set $E'\subseteq E(F_x)$ such that every vertex of $M$ is incident with two edges of $E'$, every vertex of $G\setminus M$ with one, and the $2$-factor consisting of $E'$ together with the $a$-edges at $G\setminus M$ has only even cycles.  Here $M$ is the set of vertices whose $a$-edge is not used by the even $2$-factor; there seems to be no useful algebraic characterisation of this condition.

\subsection{Normal subgroups containing the involution}
\begin{lemma}\label{lem:a-in-N}
Let $X=\X{G}{a}{x}$ with $\ord(x)$ odd.  If $G$ has a proper normal subgroup $N$ with $a\in N$ (equivalently, if there is a surjective homomorphism $\varphi\colon G\to\ZZ_m$, $m\ge3$, with $\varphi(a)=0$), then $X$ is $3$-edge-colourable.
\end{lemma}
\begin{proof}
$G/N=\langle xN\rangle$ is cyclic of order $m\mid\ord(x)$, and $m\ge3$ because $N\ne G$; let $\varphi\colon G\to\ZZ_m$ be the quotient map, so $\varphi(a)=0$ and $\varphi(x)=1$.  Choose a proper vertex $3$-colouring $c\colon\ZZ_m\to\{1,2,3\}$ of the cycle $C_m$ (it exists for every $m\ge3$).  Colour the $x$-edge $\{g,gx\}$ with $c(\varphi(g))$ and the $a$-edge $\{g,ga\}$ with the colour in $\{1,2,3\}\setminus\{c(\varphi(g)-1),c(\varphi(g))\}$; the latter is well defined because $\varphi(ga)=\varphi(g)$.  At a vertex $g$ with $\varphi(g)=i$ the two $x$-edges $\{gx^{-1},g\}$ and $\{g,gx\}$ have the distinct colours $c(i-1)$ and $c(i)$, and the $a$-edge has the third colour.
\end{proof}

Note that $m$ is automatically odd.  The lemma says that the ``degenerate quotient'' in which the $a$-edges collapse to semi-edges is never an obstacle.

\subsection{Small girth}\label{sec:girth}
A cycle of length $\ell$ through the identity of $X=\X{G}{a}{x}$ corresponds to a cyclically reduced word of length $\ell$ in $a,x,x^{-1}$ (no factor $aa$, $xx^{-1}$ or $x^{-1}x$, also cyclically) representing $1$ in $G$.

\begin{proposition}\label{prop:girth}
Let $X=\X{G}{a}{x}$ with $s=\ord(x)$ odd, $s\ge3$, and $G$ not cyclic.
\begin{enumerate}[label=\textup{(\alph*)}]
\item $X$ has girth $3$ if and only if $s=3$.
\item If $X$ has girth $4$, then $\langle x\rangle\trianglelefteq G$ and $G\cong D_{2s}$ is dihedral.
\item If $X$ has girth $5$, then either $s=5$ or $\langle x\rangle\trianglelefteq G$, so that $G=\langle x\rangle\rtimes\langle a\rangle$ is metacyclic.
\end{enumerate}
In particular Conjecture~\ref{conj:AZ} holds for all cubic Cayley graphs of girth $4$, and for those of girth $3$ (resp.\ $5$) unless $s=3$ (resp.\ $s=5$).
\end{proposition}
\begin{proof}
A cyclically reduced word in $a,x^{\pm1}$ of length $\ell\le5$ has one of the forms $x^{\pm\ell}$; $ax^{k}$ with $|k|=\ell-1$; or (for $\ell\in\{4,5\}$) $ax^{i}ax^{j}$ with $i,j\ne0$ and $|i|+|j|=\ell-2$.  A relation $x^{\pm\ell}=1$ with $\ell\le5$ and $s$ odd forces $s=\ell\in\{3,5\}$.  A relation $ax^{k}=1$ gives $a\in\langle x\rangle$ and $G$ cyclic, which is excluded.  A relation $ax^{i}ax^{j}=1$ gives $ax^{i}a=x^{-j}$ with $|i|,|j|\in\{1,2,3\}$ and $|i|+|j|\le3$; since $s$ is odd, $x^{i}$ and $x^{-j}$ generate $\langle x\rangle$ unless $s=3$ and $3\mid i$ or $3\mid j$, which would give girth $3$.  Hence $a$ normalises $\langle x\rangle$, $\langle x\rangle\trianglelefteq G$, and $G=\langle x\rangle\rtimes\langle a\rangle$.  For $\ell=4$ the only such relations are $(ax)^2=1$ (giving $D_{2s}$) and $axax^{-1}=1$ (giving $G$ abelian, hence cyclic).  The final assertion follows from Theorem~\ref{thm:ALZ} below, since metacyclic groups are solvable.
\end{proof}

Combined with Theorem~\ref{thm:cyclic}, this shows that in a proof of Conjecture~\ref{conj:AZ} by ``girth'' the cases requiring an idea are girth $3$ with $s=3$, girth $5$ with $s=5$, girth $6$, and girth at least $7$; a proof of Jaeger's conjecture on cyclically $7$-edge-connected snarks would settle the last case.  This parallels the observation of Aboomahigir and Nedela \cite{AboomahigirNedela19} that, modulo Thomassen's conjecture that cyclically $7$-connected cubic graphs other than the Coxeter graph are hamiltonian, the hamiltonicity of cubic Cayley graphs reduces to girth at most $6$.

\subsection{Generators of order three: arc-regular cubic graphs}\label{sec:arcregular}
The case $s=3$ has a pleasant reformulation.  Recall that the \emph{truncation} $\Tr(Y)$ of a cubic graph $Y$ is obtained by replacing every vertex by a triangle; its vertices are the arcs (ordered pairs of adjacent vertices) of $Y$, two arcs being adjacent if they have the same initial vertex or are inverse to each other.

\begin{lemma}\label{lem:truncation}
A cubic multigraph $Y$ is $3$-edge-colourable if and only if $\Tr(Y)$ is.
\end{lemma}
\begin{proof}
In a $3$-edge-colouring of $\Tr(Y)$ the three edges of a triangle receive three distinct colours, hence so do the three edges leaving the triangle; restricting the colouring to the edges of $\Tr(Y)$ that come from edges of $Y$ gives a $3$-edge-colouring of $Y$.  Conversely, given a $3$-edge-colouring of $Y$, colour each triangle edge of $\Tr(Y)$ with the colour of the edge of $Y$ opposite to it in the triangle.
\end{proof}

A group $G\le\Aut(Y)$ is \emph{arc-regular} if it acts regularly (transitively with trivial stabilisers) on the arcs of $Y$.

\begin{proposition}\label{prop:arcregular}
Let $G=\langle a,y\rangle$ with $a^2=y^3=1$ and $|G|>6$.  Let $Y$ be the coset graph with vertex set $\{\langle y\rangle g: g\in G\}$ in which $\langle y\rangle g$ is adjacent to $\langle y\rangle a y^{i}g$ for $i=0,1,2$.  Then $Y$ is a connected simple cubic graph on $|G|/3$ vertices, $G$ acts on it by right multiplication as an arc-regular group of automorphisms, and $\X{G}{a}{y}\cong\Tr(Y)$.  Conversely, if $Y$ is a connected cubic graph with an arc-regular group $G$ of automorphisms, then $G=\langle a,y\rangle$ for an involution $a$ and an element $y$ of order $3$, and $\Tr(Y)\cong\X{G}{a}{y}$.
\end{proposition}
\begin{proof}
Put $H=\langle y\rangle$.  The three neighbours $Ha$, $Hay$, $Hay^2$ of the vertex $H$ are distinct: $Ha=Hay^{i}$ would give $ay^{i}a^{-1}\in H$, so that $a$ normalises $H$ and $|G|\le 6$.  Right multiplication by $G$ preserves adjacency and is transitive on vertices; the stabiliser $H$ of the vertex $H$ permutes its three neighbours cyclically ($Hay^{i}\cdot y=Hay^{i+1}$), so $G$ is transitive on arcs, and $|G|=3\,|V(Y)|$ equals the number of arcs, so the action on arcs is regular.  Identify the arc $(Hg,Hag)$ with $g\in G$; arcs with the same initial vertex correspond to $g$ and $y^{\pm1}g$, and the inverse of the arc $g$ is the arc $ag$.  So $\Tr(Y)$ is the Cayley graph of $G$ with respect to $\{a,y,y^{-1}\}$ and \emph{left} multiplication, which is isomorphic to $\X{G}{a}{y}$ via $g\mapsto g^{-1}$.

Conversely, if $G$ is arc-regular on $Y$, the stabiliser $G_v$ of a vertex $v$ acts regularly on the three arcs leaving $v$, so $G_v=\langle y\rangle\cong\ZZ_3$; some $a\in G$ maps the arc $(v,w)$ to $(w,v)$, and $a^2$ fixes $(v,w)$, so $a$ is an involution.  Connectedness of $Y$ gives $G=\langle G_v,a\rangle$, and identifying arcs with group elements as above gives $\Tr(Y)\cong\X{G}{a}{y}$.
\end{proof}

\begin{corollary}\label{cor:arcregular}
Conjecture~\ref{conj:AZ} for generating pairs $(a,x)$ with $\ord(x)=3$ is equivalent to the statement that every connected cubic graph admitting an arc-regular group of automorphisms is $3$-edge-colourable.
\end{corollary}

Conder's census \cite{Conder11} of all connected cubic arc-transitive graphs on at most $10\,000$ vertices records that every one of them other than the Petersen graph and the Coxeter graph has a Hamilton cycle.  Since the Coxeter graph is $3$-edge-colourable, Lemma~\ref{lem:truncation} and Proposition~\ref{prop:arcregular} give:

\begin{corollary}\label{cor:conder}
Every connected cubic graph on at most $10\,000$ vertices that admits an arc-regular group of automorphisms is $3$-edge-colourable.  Equivalently, Conjecture~\ref{conj:AZ} holds for every $\X{G}{a}{x}$ with $\ord(x)=3$ and $|G|\le 30\,000$.
\end{corollary}

The Petersen graph is arc-transitive with automorphism group $S_5$, but has no arc-regular subgroup ($S_5$ has no subgroup of order $30$), in accordance with the conjecture; the same applies to its truncation, which is a vertex-transitive cubic graph that is not $3$-edge-colourable but not a Cayley graph.  Cubic arc-transitive graphs have been classified into seventeen ``action types'' by Conder and Nedela \cite{ConderNedela09} according to which of the Djokovi\'c--Miller groups \cite{DjokovicMiller80} occur as subgroups of the automorphism group; the ones containing an arc-regular (``$1$-regular'') subgroup are exactly the graphs in Corollary~\ref{cor:arcregular}.  Since every finite non-abelian simple group, apart from the symplectic groups $\operatorname{PSp}(4,2^f)$, $\operatorname{PSp}(4,3^f)$ and finitely many others (among them $A_6$, $A_7$, $A_8$, $\PSL(3,4)$, $\operatorname{PSU}(3,3)$, $\operatorname{PSU}(3,5)$, $\operatorname{PSU}(4,3)$, $\operatorname{PSU}(5,2)$ and the sporadic groups $M_{11}$, $M_{22}$, $M_{23}$, $McL$), is generated by an involution and an element of order $3$ (Liebeck and Shalev \cite{LiebeckShalev96}, Woldar \cite{Woldar89}; see \cite{Pellegrini15} for the classical groups of small dimension), the case $s=3$ alone already concerns a cubic Cayley graph on almost every simple group.  (In our computations, Table~\ref{tab:computation}, the groups $A_6$, $A_7$, $A_8$, $\PSL(3,4)$, $\operatorname{PSU}(3,3)$ and $M_{11}$ indeed have no generating pair with $\ord(x)=3$.)

\subsection{Normal subgroups of index two: bi-Cayley graphs}\label{sec:bicayley}
Suppose that $X=\X{G}{a}{x}$ with $s=\ord(x)$ odd and that $T\trianglelefteq G$ has index $2$.  Then $x\in T$ (its image in $G/T\cong\ZZ_2$ has odd order) and $a\notin T$, so $G=T\rtimes\langle a\rangle$ and conjugation by $a$ induces an involutory automorphism $\alpha$ of $T$ with $T=\langle x,\alpha(x)\rangle$.  The vertex set of $X$ is $T\,\sqcup\,Ta$, the $x$-edges inside $T$ are the cycles of $x$, the $x$-edges inside $Ta$ are, after the identification $ta\leftrightarrow t$, the cycles of $\alpha(x)$ (because $ta\cdot x=t\,\alpha(x)\,a$), and the $a$-edges form the identity matching $t\leftrightarrow ta$.  Thus
\begin{equation}\label{eq:bicayley}
X\;\cong\; B(T;x,\alpha(x)),
\end{equation}
where, for a group $T$ and elements $u,v\in T$, the graph $B(T;u,v)$ has vertex set $T\times\{0,1\}$ and edges $(t,0)\sim(tu,0)$, $(t,1)\sim(tv,1)$ and $(t,0)\sim(t,1)$ for all $t\in T$.
This is a \emph{bi-Cayley graph} of $T$, a natural generalisation of the generalised Petersen graphs $GP(n,k)=B(\ZZ_n;1,k)$.  Conversely, for any group $T$, any $x\in T$ of odd order and any involutory automorphism $\alpha$ of $T$ with $\langle x,\alpha(x)\rangle=T$, the graph $B(T;x,\alpha(x))$ is the cubic Cayley graph $\X{T\rtimes\langle\alpha\rangle}{\alpha}{x}$.

The order-three case of this construction is always colourable, without any assumption on $T$.

\begin{proposition}\label{prop:index2-three}
Let $G=\langle a,x\rangle$ with $a^2=x^3=1$.  If $G$ has a normal subgroup of index two, then $\X{G}{a}{x}$ is $3$-edge-colourable.
\end{proposition}
\begin{proof}
Let $T\trianglelefteq G$ have index two.  The image of $x$ in $G/T\cong\ZZ_2$ has odd order, so $x\in T$; since $a$ and $x$ generate $G$, we have $a\notin T$.  Contract every $x$-triangle of $X=\X{G}{a}{x}$ to a vertex, retaining the $a$-edges.  The resulting graph $Y$ is a cubic multigraph and $X\cong\Tr(Y)$.  Every contracted triangle lies wholly in either $T$ or $Ta$, while every $a$-edge joins these two sets.  Thus $Y$ is bipartite.  By K\H{o}nig's edge-colouring theorem $Y$ is $3$-edge-colourable, and Lemma~\ref{lem:truncation} now gives a $3$-edge-colouring of $X$.
\end{proof}

\begin{corollary}\label{cor:three-simple}
If a Cayley snark $\X{G}{a}{x}$ with $\ord(x)=3$ exists, then every one of smallest order among such graphs has $G$ non-abelian simple.
\end{corollary}
\begin{proof}
Choose such a graph with $|G|$ minimum and suppose that $1\ne N\trianglelefteq G$ is proper.  If $a\in N$, Lemma~\ref{lem:a-in-N} colours $X$.  If $a,x\notin N$, then $xN$ still has order three, and Lemma~\ref{lem:cover} makes $X$ a cover of a smaller Cayley graph; that quotient must also be a snark, contrary to minimality.  The only remaining possibility is $x\in N$ and $a\notin N$, in which case $G/N=\langle aN\rangle\cong\ZZ_2$ and Proposition~\ref{prop:index2-three} again colours $X$.  Hence $G$ is simple; it is non-abelian by Proposition~\ref{prop:trivial}.
\end{proof}

The generalised Petersen graph $GP(n,k)$ is a Cayley graph if and only if $k^2\equiv1\pmod n$ (Nedela and \v{S}koviera \cite{NS95b}), which is exactly the condition that multiplication by $k$ be an involutory automorphism of $\ZZ_n$; and Castagna and Prins \cite{CastagnaPrins72} proved Watkins' conjecture that every generalised Petersen graph other than the Petersen graph itself is $3$-edge-colourable.  (The Petersen graph is $GP(5,2)$ with $2^2\equiv-1$, not $+1$, modulo $5$.)  So Conjecture~\ref{conj:AZ} holds for all generalised Petersen graphs, and the results of \S\ref{sec:minimal} show that the whole conjecture reduces to graphs of the form \eqref{eq:bicayley} with $T$ simple or $T\cong S\times S$, together with Cayley graphs on simple groups.

\section{Quotients by arbitrary subgroups, and a reformulation over $\mathbb{Z}_3$}\label{sec:new}

This section contains two elementary observations which we have not found in the literature in this form.  The first extends the quotient lemma (Lemma~\ref{lem:cover}) from normal subgroups to arbitrary subgroups, at the price of allowing semi-edges in the quotient; it is the tool behind the computations of Section~\ref{sec:computation}.  The second describes the $3$-edge-colourings of $\X{G}{a}{x}$ as the $a$-invariant functions $G\to\ZZ_3$ satisfying a local condition along the $x$-cycles.

\subsection{Pregraph quotients}\label{sec:pregraph}
A \emph{pregraph} is a graph in which, besides ordinary edges, \emph{semi-edges} (edges with only one end) and loops are allowed; each vertex of a cubic pregraph is incident with three \emph{darts}, an ordinary edge contributing a dart at each of its ends and a semi-edge one dart.  A $3$-edge-colouring of a pregraph colours the ordinary edges and the semi-edges so that the three darts at every vertex carry three different colours.

Let $X=\X{G}{a}{x}$ with $s=\ord(x)$ odd and let $H\le G$ be \emph{any} subgroup.  The left action of $H$ on $V(X)=G$ is free, and so is its action on the darts $(g,gs)$, $s\in\{a,x,x^{-1}\}$.  The quotient pregraph $H\backslash X$ has vertex set $H\backslash G=\{Hg\}$ and one dart for each $H$-orbit of darts; the darts $(g,gs)$ and $(gs,g)$ lie in the same orbit if and only if $hg=gs$ and $hgs=g$ for some $h\in H$, i.e.\ $s^2=1$ and $gsg^{-1}\in H$, which happens only for $s=a$.  Hence the $a$-edge at $Hg$ is a semi-edge exactly when $gag^{-1}\in H$, and the $x$-edge at $Hg$ is a loop exactly when $gxg^{-1}\in H$.

\begin{lemma}[generalised quotient lemma]\label{lem:pregraph}
Let $X=\X{G}{a}{x}$ with $\ord(x)$ odd and let $H\le G$ be a subgroup containing no conjugate of $x$.  Then $\bar X=H\backslash X$ is a loopless cubic pregraph on $|G:H|$ vertices with exactly $|H\cap a^G|\cdot|C_G(a)|/|H|$ semi-edges, and $X\to\bar X$, $g\mapsto Hg$, is a covering of pregraphs.  If $\bar X$ is $3$-edge-colourable, then so is $X$; more precisely, the $3$-edge-colourings of $\bar X$ correspond bijectively to the $3$-edge-colourings of $X$ that are invariant under left multiplication by $H$.
\end{lemma}
\begin{proof}
The three darts at $g$ are mapped to the three darts at $Hg$, which are distinct because $\bar X$ has no loops; so the projection is a covering.  Given a colouring of $\bar X$, colour every dart of $X$ with the colour of its image; the two darts of an ordinary edge $\{g,ga\}$ with $gag^{-1}\notin H$ map to the two darts of an ordinary edge of $\bar X$ and receive the same colour, and if $gag^{-1}\in H$ they map to the same semi-edge.  Thus the colouring is a well defined $3$-edge-colouring of $X$, obviously $H$-invariant, and conversely an $H$-invariant colouring of $X$ is constant on $H$-orbits of darts and descends to $\bar X$.  The count of semi-edges is the number of cosets $Hg$ with $gag^{-1}\in H$, i.e.\ the number of elements $g$ with $gag^{-1}\in H\cap a^G$, which is $|H\cap a^G|\,|C_G(a)|$, divided by $|H|$.
\end{proof}

Lemma~\ref{lem:cover} is the case $H\trianglelefteq G$ with $a\notin H$, and Lemma~\ref{lem:lift} in the form used by the program of \S\ref{sec:computation} is the case of odd $|H|$ (no semi-edges).  The point of Lemma~\ref{lem:pregraph} is that $H$ may now be large: any subgroup whose order is not divisible by $\ord(x)$, the centraliser $C_G(a)$ whenever no conjugate of $x$ centralises $a$ (the quotient is then the pregraph on the conjugacy class $a^G$ in which $b$ is joined to $b^a$, $b^x$ and $b^{x^{-1}}$), a point stabiliser of any permutation representation in which $x$ has no fixed points, and so on.  Semi-edges carry a parity constraint which, in practice, is the main reason why the largest admissible subgroup need not work.

\begin{lemma}[parity]\label{lem:parity}
Let $\bar X$ be a cubic pregraph with $n$ vertices and $\sigma$ semi-edges.  If $\bar X$ is $3$-edge-colourable then the semi-edges fall into three colour classes each of size $\equiv n\pmod 2$; in particular $\sigma\equiv n\pmod 2$, and $\sigma\ge3$ if $n$ is odd.
\end{lemma}
\begin{proof}
Each colour class contains $e$ ordinary edges and $\sigma'$ semi-edges with $2e+\sigma'=n$.
\end{proof}

\begin{example}\label{ex:projline}
Let $G=\PSL(2,q)$ act on the projective line $\Omega=\GF(q)\cup\{\infty\}$, and let $x$ lie in a non-split torus, so that $x$ has no fixed point on $\Omega$ and $H=G_\infty$ (a Borel subgroup) is admissible in Lemma~\ref{lem:pregraph}; the quotient is the Schreier graph of $\langle a,x\rangle$ on the $q+1$ points.  If $q$ is even, an involution $a$ fixes exactly one point, so $\bar X$ has $q+1$ (odd) vertices and one semi-edge, and Lemma~\ref{lem:parity} shows that $\bar X$ is never $3$-edge-colourable---although $X$ is, for $q\in\{4,8,16,32,64\}$ (Table~\ref{tab:computation}).  If $q\equiv 3\pmod 4$ the involutions are fixed-point-free on $\Omega$ and if $q\equiv1\pmod4$ they have two fixed points, so the parity condition holds; for these $q$ we found by computer that the $(q+1)$-vertex quotient is $3$-edge-colourable for most but not all generating pairs (for instance for all $192$ generating pairs with $\ord(x)=13$ in $\PSL(2,25)$ or $\ord(x)=15$ in $\PSL(2,29)$, but for only $11$ of the $13$ generating pairs with $\ord(x)=\ord(ax)=7$ in $\PSL(2,13)$).
\end{example}

\subsection{Hamiltonian and cyclewise quotients}\label{sec:hamquot}
The quotient $H\backslash X$ is itself ``hamiltonian'' when $\langle x\rangle$ acts transitively on $H\backslash G$, i.e.\ when $G=H\langle x\rangle$ is an exact factorisation; the $x$-cycle through $H$ then passes through all $|G:H|=\ord(x)$ vertices.  The following sufficient condition also applies when the $x$-edges form several cycles.

\begin{lemma}[transition criterion]\label{lem:hampregraph}
Let $Y$ be a loopless cubic pregraph with a spanning $2$-factor $F$ all of whose cycles are odd.  Contract the cycles of $F$ to obtain a pregraph $Q$, retaining the complementary darts in their cyclic order around each contracted vertex.  Suppose that $Q$ has a spanning $2$-regular subpregraph $R$ which uses two consecutive darts at every vertex and whose circuit components all have even length.  Then $Y$ is $3$-edge-colourable.
\end{lemma}
\begin{proof}
At each vertex of $Q$, the two darts selected by $R$ correspond to complementary darts of $Y$ at adjacent vertices $u,u'$ of a cycle $C$ of $F$.  Delete the edge $uu'$ from $C$ and retain the two selected darts.  Doing this on every cycle of $F$ produces a spanning $2$-regular subpregraph $Z$ of $Y$.  A circuit component of $R$ expands to a circuit of $Z$ whose length has the same parity: every visit to $C$ replaces a vertex by the path $C-uu'$ of even length.  Hence every circuit of $Z$ is even.  Every other component of $R$, and therefore of $Z$, is a path capped by two semi-edges.  Colour the circuits of $Z$ alternately $0,1$, do the same along each capped path (including its two semi-edges), and colour every dart outside $Z$ with $2$.  Every vertex sees all three colours.
\end{proof}

\begin{corollary}\label{cor:cyclewise-quotient}
Let $X=\X{G}{a}{x}$ with $\ord(x)$ odd, and let $H\le G$ contain no conjugate of $x$.  Suppose that every cycle of the permutation $Hg\mapsto Hgx$ of $H\backslash G$ contains adjacent cosets $Hg,Hgx$ for which $gag^{-1},gxax^{-1}g^{-1}\in H$.  Then $X$ is $3$-edge-colourable.
\end{corollary}
\begin{proof}
By Lemma~\ref{lem:pregraph}, $H\backslash X$ is a loopless cubic pregraph.  Its $x$-edges form a spanning $2$-factor of odd cycles, and the displayed membership conditions say precisely that the complementary darts at $Hg$ and $Hgx$ are semi-edges.  In the contracted pregraph $Q$ of Lemma~\ref{lem:hampregraph}, choose those two adjacent semi-edges at every vertex; their union is a spanning $2$-regular subpregraph with no circuit components.  Lemma~\ref{lem:hampregraph} colours the quotient, and Lemma~\ref{lem:pregraph} lifts the colouring to $X$.
\end{proof}

\begin{proposition}\label{prop:exact}
Let $X=\X{G}{a}{x}$ with $\ord(x)$ odd.  Suppose $G=H\langle x\rangle$ with $H\cap\langle x\rangle=1$ for a subgroup $H$ containing $a$ and $xax^{-1}$.  Then $X$ is $3$-edge-colourable.
\end{proposition}
\begin{proof}
Since $\langle x\rangle$ acts regularly on $H\backslash G$, $x$ has no fixed point there, so $H$ contains no conjugate of $x$ and Lemma~\ref{lem:pregraph} applies; the pregraph $H\backslash X$ has the Hamilton cycle $H,Hx,Hx^2,\dots$, and its adjacent vertices $H$ and $Hx$ carry semi-edges because $a\in H$ and $xax^{-1}\in H$.  Apply Lemma~\ref{lem:hampregraph}.
\end{proof}

For instance, if $n$ is odd, $x$ is an $n$-cycle in $S_n$ and $a$ is an involution fixing two cyclically consecutive points of $x$ (say $a=(1\,2)$, or any involution with at least $(n+1)/2$ fixed points), then Proposition~\ref{prop:exact} with $H=S_{n-1}$ shows that $\X{\langle a,x\rangle}{a}{x}$ is $3$-edge-colourable; for $a=(1\,2)$ this is a two-line proof of the colouring consequence of the Compton--Williamson theorem quoted in \S\ref{sec:hamilton}.  The adjacency hypothesis cannot simply be replaced by ``at least three semi-edges'' (the minimum allowed by Lemma~\ref{lem:parity}): the Petersen graph minus a vertex is a $9$-cycle with three chords and three vertices of valency two, i.e.\ the pregraph $S_8\backslash X$ for $x=(0\,1\,\cdots\,8)$ and $a=(1\,5)(2\,7)(4\,8)$, and it is not $3$-edge-colourable, although $\langle a,x\rangle$ (a solvable group of order $162$) and $X$ are perfectly harmless.  This example illustrates the main difficulty of the quotient method: a small ``natural'' quotient of a colourable Cayley graph can be a Petersen-like snark, so a proof of Conjecture~\ref{conj:AZ} along these lines must be able to choose the subgroup $H$ adaptively.

\subsection{The colourings as functions $G\to\ZZ_3$}\label{sec:Z3}
Let $X=\X{G}{a}{x}$ with $s=\ord(x)$ odd, and identify the three colours with the elements of $\ZZ_3$, so that the three colours at a vertex sum to $0$.  Given a $3$-edge-colouring $\kappa$ of $X$, let $t(g)\in\ZZ_3$ be the colour of the $a$-edge at $g$, i.e.\ the colour missing from the two $x$-edges at $g$.  Then $t(ga)=t(g)$, and $\kappa$ is determined by $t$: on the $x$-cycle through $g$, the colour of the edge $\{gx^{i},gx^{i+1}\}$ is $-\bigl(t(gx^{i})+\kappa(\{gx^{i-1},gx^i\})\bigr)$, and the following lemma shows that $t$ determines $\kappa$ on the cycle without any initial choice.  We say that a function $t\colon\ZZ_s\to\ZZ_3$ is \emph{admissible} if it is the sequence of colours of the pendant edges in some proper $3$-edge-colouring of the cycle $C_s$ (with a pendant edge attached at each vertex), equivalently the sequence of colours missing at the vertices of $C_s$.

\begin{lemma}\label{lem:admissible}
Let $s\ge3$ be odd and $t\colon\ZZ_s\to\ZZ_3$.  Call $i\in\ZZ_s$ a \emph{change} if $t(i+1)\ne t(i)$, and put $\delta_i=t(i+1)-t(i)\in\{\pm1\}$ for a change $i$.  Then $t$ is admissible if and only if $t$ is not constant and, for any two cyclically consecutive changes $i$ and $j$ (so that $t$ is constant on $i+1,\dots,j$), $\delta_i=\delta_j$ holds precisely when $j-i$ is odd.  Each admissible $t$ comes from exactly one $3$-edge-colouring of $C_s$; in particular there are $2^s-2$ admissible functions.
\end{lemma}
\begin{proof}
Let $f_i\in\ZZ_3$ be the colour of the edge $\{i,i+1\}$ of $C_s$ in a proper $3$-edge-colouring, so $t(i)=-(f_{i-1}+f_i)$, and put $\epsilon_i=f_i-f_{i-1}\in\{\pm1\}$.  Then $t(i+1)-t(i)=-(\epsilon_i+\epsilon_{i+1})$, so $i$ is a change if and only if $\epsilon_i=\epsilon_{i+1}$, in which case $\delta_i=-2\epsilon_i=\epsilon_i$, while $\epsilon_{i+1}=-\epsilon_i$ at a non-change.  Between consecutive changes $i<j$ the signs $\epsilon_{i+1},\dots,\epsilon_j$ alternate, so $\delta_j=\epsilon_j=(-1)^{j-i-1}\epsilon_{i+1}=(-1)^{j-i-1}\delta_i$, which is the stated rule; and $t$ is not constant since $s$ is odd.  Conversely, let $t$ be non-constant and satisfy the rule.  Define $\epsilon_{i}=\epsilon_{i+1}=\delta_i$ at every change $i$ and let $\epsilon$ alternate between changes; the rule says exactly that this is consistent around the cycle.  Put $f_i=f_0+\epsilon_1+\dots+\epsilon_i$; this is a proper colouring of $C_s$ provided $\epsilon_1+\dots+\epsilon_s\equiv0\pmod 3$, and indeed $\sum_i\epsilon_i\equiv-2\sum_i\epsilon_i=\sum_i\bigl(t(i+1)-t(i)\bigr)=0$.  By construction $-(f_{i-1}+f_i)$ and $t(i)$ have the same differences, so they differ by a constant $c$, and replacing $f$ by $f+c$ (which changes $-(f_{i-1}+f_i)$ by $-2c=c$) makes them equal.  Since $\epsilon$ and $c$ are forced, the colouring is unique; the count follows because $C_s$ has $2^s-2$ proper $3$-edge-colourings for odd $s$.
\end{proof}

\begin{proposition}\label{prop:Z3}
Let $X=\X{G}{a}{x}$ with $s=\ord(x)$ odd.  The $3$-edge-colourings of $X$ (with colours in $\ZZ_3$) are in bijection with the functions $t\colon G\to\ZZ_3$ such that $t(ga)=t(g)$ for all $g$ and, for every $g$, the sequence $i\mapsto t(gx^i)$ on $\ZZ_s$ is admissible.
\end{proposition}
\begin{proof}
The map $\kappa\mapsto t$ described above is injective by Lemma~\ref{lem:admissible}, and conversely, given $t$, colouring each $x$-cycle by the unique colouring of Lemma~\ref{lem:admissible} and each $a$-edge $\{g,ga\}$ with $t(g)=t(ga)$ gives a proper colouring: at $g$ the two $x$-edges have the colours other than $t(g)$.
\end{proof}

Thus a $3$-edge-colouring of $X$ is the same thing as a $3$-colouring of the perfect matching $M_a$ (of the $|G|/2$ $a$-edges) such that, along every $x$-cycle, the colours of the $a$-edges met form an admissible sequence: a constraint satisfaction problem with $|G|/2$ variables and $|G|/s$ constraints of arity $s$, each variable occurring in two constraints.  The simplest admissible sequences are those with exactly three changes, all with the same $\delta$: $t$ takes the values $c,c+1,c+2$ on three consecutive arcs of odd lengths.  Lemma~\ref{lem:a-in-N} is the case in which $t$ is pulled back from a valid function on the quotient cycle $\ZZ_m=G/N$.  Two small cases deserve to be spelled out.

\begin{corollary}\label{cor:s35}
Let $X=\X{G}{a}{x}$.
\begin{enumerate}[label=\textup{(\alph*)}]
\item If $\ord(x)=3$, then $X$ is $3$-edge-colourable if and only if there is $t\colon G\to\ZZ_3$ with $t(ga)=t(g)$ which takes three different values on every triangle $\{g,gx,gx^2\}$, i.e.\ if and only if the $4$-valent graph $X/M_a$ obtained by contracting the $a$-edges has a proper $3$-vertex-colouring.  \textup{(}In the notation of \S\ref{sec:arcregular}, $X/M_a$ is the line graph of the arc-regular graph $Y$, so this is Lemma~\ref{lem:truncation} again.\textup{)}
\item If $\ord(x)=5$, then $X$ is $3$-edge-colourable if and only if there is $t\colon G\to\ZZ_3$ with $t(ga)=t(g)$ such that on every $x$-cycle $t$ is constant on three consecutive vertices and takes the two other values on the remaining two vertices.
\end{enumerate}
\end{corollary}
\begin{proof}
For $s=3$ a sequence is admissible if and only if its three values are distinct (there are $6=2^3-2$ of them).  For $s=5$ the $30=2^5-2$ admissible sequences are exactly the rotations of $(c,c,c,c',c'')$ with $\{c,c',c''\}=\ZZ_3$: these are admissible by Lemma~\ref{lem:admissible} (three changes at mutual distances $3,1,1$, all with the same $\delta$), and there are $5\cdot 6=30$ of them.
\end{proof}

There is a useful way to isolate the remaining difficulty in the index-two case.  Suppose $T\trianglelefteq G$ has index two, and contract each cycle of $F_x$ to a vertex while retaining the $a$-edges.  The resulting graph $Q$ is a bipartite $s$-regular multigraph: its two parts consist of the $x$-cycles in $T$ and in $Ta$.  The order of the vertices along each $x$-cycle gives a cyclic order to the $s$ edges incident with the corresponding vertex of $Q$.  Call a $2$-factor of $Q$ \emph{consecutive} if its two edges at every vertex are consecutive in this cyclic order.  In the terminology of Kratochv\'il and Poljak \cite{KratochvilPoljak92}, this is a compatible $2$-factor for the transition system in which the allowed pairs at each vertex form the cycle $C_s$.

\begin{proposition}[consecutive-$2$-factor criterion]\label{prop:consecutive}
Let $X=\X{G}{a}{x}$ with $s=\ord(x)$ odd, and suppose that $G$ has a normal subgroup of index two.  If the cyclically ordered bipartite multigraph $Q$ defined above has a consecutive $2$-factor, then $X$ is $3$-edge-colourable.
\end{proposition}
\begin{proof}
Every cycle of the $2$-factor is even because $Q$ is bipartite.  Colour its edges alternately $1,2$, and colour every edge outside the $2$-factor with $0$.  Transfer this colour to the corresponding $a$-edge of $X$.  Around each $x$-cycle the colours of the incident $a$-edges are a rotation of
\[
  (\underbrace{0,\ldots,0}_{s-2},1,2).
\]
This sequence is admissible: its three changes have the same difference in $\ZZ_3$ and occur at mutual cyclic distances $s-2,1,1$, all odd.  Proposition~\ref{prop:Z3} therefore supplies a $3$-edge-colouring of $X$.
\end{proof}

For $s=3$, K\H{o}nig's theorem decomposes the cubic bipartite multigraph $Q$ into three perfect matchings; the union of any two is a $2$-factor, and every pair among three cyclically ordered incident edges is consecutive.  Thus Proposition~\ref{prop:consecutive} recovers Proposition~\ref{prop:index2-three}.  For $s\ge5$, regularity and bipartiteness still guarantee a $2$-factor, but not that it is consecutive.  In the notation of \eqref{eq:bicayley}, $Q$ is the coset-incidence graph whose two parts are the left cosets of $\langle x\rangle$ and $\langle\alpha(x)\rangle$ in $T$, with one edge labelled by each $t\in T$; the open task is therefore to exploit this group-induced cyclic ordering to find a consecutive $2$-factor.

Proposition~\ref{prop:Z3} combines with Lemma~\ref{lem:pregraph}: an $H$-invariant colouring is a function $t$ on $H\backslash G$, constant on the orbits of $a$, admissible on the cycles of $x$.  For the coset space of a large subgroup this is a very small constraint satisfaction problem, and it is the form in which the SAT solver of Section~\ref{sec:computation} sees the problem.  We have not been able to turn either formulation into a proof of Conjecture~\ref{conj:AZ} for any infinite family of simple groups; the difficulty is visible already in Example~\ref{ex:projline}.

\section{Solvable groups}\label{sec:solvable}

\begin{theorem}[Alspach, Liu, Zhang \cite{ALZ96}]\label{thm:ALZ}
Every connected cubic Cayley graph on a solvable group is $3$-edge-colourable.  Equivalently, every connected Cayley graph of valency at least $2$ on a solvable group admits a nowhere-zero $4$-flow.
\end{theorem}

The proof is an induction along a normal series.  With the lemmas of \S\S\ref{sec:prelim}--\ref{sec:easy} in hand it can be organised as follows.  Let $X=\X{G}{a}{x}$ with $G$ solvable and $\ord(x)$ odd, and let $N$ be a minimal normal subgroup of $G$; it is elementary abelian.  If $a\notin N$ and $x\notin N$, Lemma~\ref{lem:cover} reduces to the smaller solvable group $G/N$.  If $a\in N$, Lemma~\ref{lem:a-in-N} applies (note $N\ne G$ as $G$ is not abelian of prime order).  Otherwise $x\in N$, $N$ has index $2$ and $G=N\rtimes\langle a\rangle$ with $N$ an elementary abelian $p$-group ($p$ odd) generated by $x$ and $x^a$; thus $N\cong\ZZ_p$ or $\ZZ_p^2$, and $X$ is either the prism over $C_p$ (if $x^a=x^{\pm1}$), or the bi-Cayley graph $B(\ZZ_p^2;(1,0),(0,1))$; the prism is hamiltonian, and the second graph is $3$-edge-colourable as well (see \cite{ALZ96} for the details in the original formulation; we also confirmed this by computer for $p\le 13$).

Theorem~\ref{thm:ALZ} covers, among others, all cubic Cayley graphs on abelian, nilpotent, supersolvable, dihedral and generalised dihedral groups, on $p$-groups, on groups of order $p^aq^b$ (Burnside) and on all groups of order less than $60$.

\subsection{Solvable vertex-transitive groups}
Poto\v{c}nik extended Theorem~\ref{thm:ALZ} from regular to transitive actions of solvable groups.

\begin{theorem}[Poto\v{c}nik \cite{Potocnik04}]\label{thm:potocnik}
Let $X$ be a connected cubic simple graph whose automorphism group contains a solvable subgroup acting transitively on the vertices.  If $X$ is not $3$-edge-colourable, then $X$ is the Petersen graph.
\end{theorem}

The Petersen graph is a genuine exception: the subgroup of $S_5$ of order $20$ (the affine group $\mathrm{AGL}(1,5)$) is solvable and acts transitively on its ten vertices.  Theorem~\ref{thm:potocnik} contains Theorem~\ref{thm:ALZ}, since a Cayley graph on a solvable group admits that group as a solvable regular, hence transitive, group of automorphisms, and the Petersen graph is not a Cayley graph.

The strategy of the proof is again inductive, using the theory of elementary abelian regular covers developed by Malni\v{c}, Maru\v{s}i\v{c} and Poto\v{c}nik \cite{MMP04,MalnicPotocnik06}.  If $G\le\Aut(X)$ is solvable and vertex-transitive, a minimal normal subgroup $N$ of $G$ is elementary abelian; if $N$ acts semiregularly on $V(X)$ then $X$ is a regular $N$-cover of the quotient $X/N$, which admits the solvable vertex-transitive group $G/N$, and Lemma~\ref{lem:lift} applies provided $X/N$ is a cubic graph.  The remaining cases (non-semiregular $N$, quotients of smaller valency) are analysed directly, and the base of the induction involves the Petersen graph: as Malni\v{c} and Poto\v{c}nik put it, ``the fact that all elementary abelian covers of the Petersen graph are $3$-edge-colourable was used \dots\ in order to prove, by induction, that every connected cubic graph admitting a vertex-transitive solvable group of automorphisms (with the sole exception of the Petersen graph itself) is $3$-edge-colourable'' \cite[\S1]{MalnicPotocnik06}.

Theorem~\ref{thm:potocnik} also yields the conjecture for many non-Cayley vertex-transitive cubic graphs, and for every cubic vertex-transitive graph whose order $n$ is such that all transitive permutation groups of degree $n$ are solvable.

\section{The structure of a minimal counterexample}\label{sec:minimal}

Nedela and \v{S}koviera proved that a smallest Cayley snark, if one exists, must come from a group that is ``almost as simple as possible''.  A group $G$ is \emph{almost simple} if $T\le G\le\Aut(T)$ for some non-abelian simple group $T$.

\begin{theorem}[Nedela--\v{S}koviera \cite{NS01}]\label{thm:NS}
Let $X=\X{G}{a}{x}$ be a Cayley snark of smallest possible order.  Then either
\begin{enumerate}[label=\textup{(\roman*)}]
\item $G$ is a non-abelian simple group; or
\item $G$ has exactly one non-trivial proper normal subgroup $T$; moreover $|G:T|=2$, $x\in T$, $a\notin T$, and $T$ is either a non-abelian simple group (so that $G$ is almost simple) or the direct product $S\times S$ of two isomorphic non-abelian simple groups, the involution $a$ interchanging the two factors (so that $G\cong S\wr\ZZ_2$).
\end{enumerate}
In particular every Cayley snark is a regular cover of a Cayley snark on a group of type \textup{(i)} or \textup{(ii)}.
\end{theorem}

The last sentence follows from the first by repeated application of Lemma~\ref{lem:cover} to a Cayley snark that is not minimal (the quotient by a normal subgroup avoiding $a$ and $x$ is again a Cayley snark, by Lemma~\ref{lem:cover}, or the normal subgroup contains $a$, which is impossible by Lemma~\ref{lem:a-in-N}, or it has index $2$).  The reader should consult \cite{NS01} for the original proof; the following short argument, based on the lemmas above and on Theorem~\ref{thm:ALZ}, recovers the statement.

\begin{proof}[Proof of Theorem~\ref{thm:NS}]
Let $X=\X{G}{a}{x}$ be a Cayley snark of smallest order; by Proposition~\ref{prop:trivial}, $s=\ord(x)$ is odd and $G$ is not abelian.  Let $1\ne N\trianglelefteq G$ be a proper normal subgroup.
\begin{itemize}[leftmargin=2em]
\item If $a\notin N$ and $x\notin N$, then by Lemma~\ref{lem:cover} $X$ covers the smaller cubic Cayley graph $\X{G/N}{\bar a}{\bar x}$, which is $3$-edge-colourable by minimality; hence so is $X$, a contradiction.
\item If $a\in N$, then $X$ is $3$-edge-colourable by Lemma~\ref{lem:a-in-N}, a contradiction.
\item Hence $x\in N$ and $a\notin N$.  Then $G/N=\langle aN\rangle$ has order $2$.
\end{itemize}
So every non-trivial proper normal subgroup of $G$ has index $2$ and contains $x$.  If $N_1\ne N_2$ were two of them, $N_1\cap N_2$ would be a normal subgroup of index $4$ containing $x\ne1$, contradicting what we just proved.  Hence $G$ has at most one non-trivial proper normal subgroup.

If it has none, $G$ is simple, and non-abelian because $|G|$ is even and divisible by $s\ge3$.

Otherwise let $T$ be the unique non-trivial proper normal subgroup.  It is a minimal normal subgroup, hence characteristically simple: $T\cong S^k$ for a simple group $S$.  If $S$ is abelian then $G=T\rtimes\langle a\rangle$ is solvable, contradicting Theorem~\ref{thm:ALZ}.  So $S$ is non-abelian simple, and the $k$ direct factors of $T$ are the minimal normal subgroups of $T$, which $a$ permutes by conjugation.  The product of the factors in one $\langle a\rangle$-orbit is normalised by $T$ and by $a$, hence normal in $G$, hence equal to $T$; so $a$ acts transitively on the $k$ factors and $k\le2$.  If $k=1$, $T$ is non-abelian simple and $G$ is almost simple.  If $k=2$, $a$ interchanges the two factors; after replacing the second factor's coordinates by their image under a suitable automorphism of $S$ we may assume that $a$ acts by $(u,v)\mapsto(v,u)$, i.e.\ $G\cong S\wr\ZZ_2$.
\end{proof}

\begin{remark}
(1) In case (ii) the graph $X$ is the bi-Cayley graph $B(T;x,\alpha(x))$ of \eqref{eq:bicayley}, where $\alpha$ is the outer involutory automorphism of $T$ induced by $a$ (it is not inner: if $a$ induced an inner automorphism $\iota_t$, then $at^{-1}$ would be a central involution of $G$, generating a second normal subgroup).  Examples of admissible pairs $(G,T)$: $(S_n,A_n)$ with $a$ an odd permutation; $(\PGL(2,q),\PSL(2,q))$ for odd $q$; $(\PSigmaL(2,q^2),\PSL(2,q^2))$; $(M_{10},A_6)$; $(S\wr\ZZ_2,S\times S)$ with, for instance, $S=A_5$ and $x=\bigl((123),(12345)\bigr)$.

(2) Groups of type (i) or (ii) are precisely the groups $G=\langle a,x\rangle$ whose \emph{socle} is a non-abelian minimal normal subgroup $T$ with $G/T$ of order at most $2$; this explains the title of \cite{NS01}.  Note that the theorem does not say that every Cayley snark lives on such a group, only that a \emph{smallest} one does; a Cayley snark on an arbitrary group covers one of these.  When $\ord(x)=3$, Proposition~\ref{prop:index2-three} eliminates type~(ii), so Corollary~\ref{cor:three-simple} reduces that case to non-abelian simple groups alone.

(3) The argument shows more generally that Conjecture~\ref{conj:AZ} holds for every group $G$ that has a non-trivial normal subgroup $N$ with $a\in N$ or with $a,x\notin N$ and the conjecture known for $G/N$.  For instance, it holds for every group of the form $G=N\rtimes\ZZ_m$ with $m$ odd and $N$ containing $a$, and for all groups $G$ with a solvable non-trivial normal subgroup $N$ such that $G/N$ satisfies the conjecture (by induction, if $a,x\notin N$; by Lemma~\ref{lem:a-in-N} if $a\in N$; and if $x\in N$ then $|G:N|=2$ and $G$ is solvable).  In particular, the conjecture holds for all groups with a non-trivial solvable radical whose quotient by the radical satisfies it.
\end{remark}

\section{The hamiltonicity route and Cayley maps}\label{sec:hamilton}

As noted in the introduction, a Hamilton cycle in a cubic Cayley graph yields a $3$-edge-colouring.  The literature on Hamilton cycles in Cayley graphs is vast (see the surveys \cite{CurranGallian96,KutnarMarusic09,PakRadoicic09}); we mention only the results that bear directly on cubic Cayley graphs on non-solvable groups, since for solvable groups Theorem~\ref{thm:ALZ} is already available.

\subsection{$(2,s,3)$-Cayley graphs}
Glover and Maru\v{s}i\v{c} \cite{GM07} initiated the study of cubic Cayley graphs $\X{G}{a}{x}$ on groups with a \emph{$(2,s,3)$-presentation}
\[
G=\langle a,x \mid a^2=x^s=(ax)^3=1,\dots\rangle ,
\]
i.e.\ generated by an involution $a$ and an element $x$ of order $s$ whose product has order $3$.  Such a graph has a natural embedding as a \emph{Cayley map} in an orientable surface whose faces are the $s$-gons bounded by the $x$-cycles and the hexagons bounded by the closed walks $(ax)^3$; a Hamilton cycle is sought as the boundary of a tree of faces (a ``tree of faces'' in the Cayley map), using classical results on cyclic stability of cubic graphs.

\begin{theorem}[Glover--Maru\v{s}i\v{c} \cite{GM07}; Glover--Kutnar--Maru\v{s}i\v{c} \cite{GKM09}; Glover--Kutnar--Malni\v{c}--Maru\v{s}i\v{c} \cite{GKMM12}]\label{thm:2s3}
Let $X=\X{G}{a}{x}$ where $G$ has a $(2,s,3)$-presentation.  Then $X$ has a Hamilton path, and it has a Hamilton cycle whenever $|G|\equiv2\pmod4$, whenever $s\equiv0\pmod4$, and whenever $s$ is odd.  The only open case is $|G|\equiv0\pmod4$ with $s\equiv2\pmod4$.
\end{theorem}

\begin{corollary}\label{cor:2s3}
Every $(2,s,3)$-Cayley graph is $3$-edge-colourable.
\end{corollary}
\begin{proof}
If $s$ is even, Proposition~\ref{prop:trivial}(b) applies; if $s$ is odd, Theorem~\ref{thm:2s3} provides a Hamilton cycle.
\end{proof}

Note that a $(2,s,3)$-generated group is the same thing as a $(2,3)$-generated group (put $y=ax$), but the Cayley graph $\X{G}{a}{x}$ of Corollary~\ref{cor:2s3} is different from the graph $\X{G}{a}{y}$ of \S\ref{sec:arcregular}: in the former the element of order $3$ is the \emph{product} of the two generators.  The two families overlap exactly when $s=3$.

\subsection{Cayley maps and independent sets}
Hujdurovi\'c, Kutnar and Maru\v{s}i\v{c} \cite{HKM13,HKM14} pursued the Cayley-map approach with Conjecture~\ref{conj:AZ} rather than hamiltonicity as the target.  In \cite{HKM13} they show that a connected bicirculant of prime valency admitting a suitable automorphism is \emph{near-bipartite} (its vertex set has an independent subset whose removal leaves a bipartite graph), and combine this with the theory of Cayley maps to obtain partial results on the non-existence of Cayley snarks; \cite{HKM14} develops this further, ``combining the theory of Cayley maps with the existence of certain kinds of independent sets of vertices in arc-transitive graphs'' to obtain ``new partial results suggesting promising future research directions'' for the conjecture.  The underlying idea is that, in $\X{G}{a}{x}$ with $s$ odd, the $x$-cycles are the vertices of an $s$-valent $G$-arc-transitive quotient graph (the coset graph of $\langle x\rangle$), and a suitable independent set in this quotient allows one to modify the odd $2$-factor $F_x$ into an even one.

\subsection{Other hamiltonicity results}
Kutnar, Maru\v{s}i\v{c}, Morris, Morris and \v{S}parl \cite{KMMMS12} proved that every connected Cayley graph whose order is $kp$ with $k<32$, $k\ne 24$ ($p$ prime), or $kpq$ with $k\le5$, or $pqr$, or $kp^2$ with $k\le 4$, or $kp^3$ with $k\le2$ ($p,q,r$ distinct primes) is hamiltonian, and Morris and Wilk \cite{MorrisWilk20} extended the first case to all $k<48$; in particular every connected Cayley graph of order less than $48$ is hamiltonian, and so are, for instance, the cubic Cayley graphs of order $60$ on $A_5$ and of order $120$ on $S_5$ and $\mathrm{SL}(2,5)$.  All of these are, however, well inside the range of the census discussed next.

A different kind of result concerns a single, very natural family: for $\sigma=(1\,2\,\cdots\,n)$ and $\tau=(1\,2)$ the graph $\X{S_n}{\tau}{\sigma}$ is the undirected ``sigma-tau graph''.  Compton and Williamson \cite{ComptonWilliamson93} showed that it is hamiltonian for every $n\ge3$, and Sawada and Williams \cite{SawadaWilliams20} solved the much harder directed version (Nijenhuis--Wilf's sigma-tau problem): the Cayley \emph{digraph} on $\{\sigma,\tau\}$ has a Hamilton cycle if and only if $n$ is odd.  For odd $n$ this is a bi-Cayley graph of the form \eqref{eq:bicayley} on $T=A_n$ with $\alpha$ conjugation by a transposition, so Conjecture~\ref{conj:AZ} holds for these graphs; the case of a general odd-order $x\in A_n$ is open.

\section{Computational evidence}\label{sec:computation}

\subsection{The census of cubic vertex-transitive graphs}
Poto\v{c}nik, Spiga and Verret \cite{PSV13} constructed all connected cubic vertex-transitive graphs on at most $1280$ vertices: there are $111\,360$ of them, of which $109\,926$ are Cayley graphs ($97\,687$ of these are graphical regular representations, i.e.\ their full automorphism group acts regularly).  They report:

\begin{quote}
``We tested all the graphs in the census for Hamilton cycles and found at least one such cycle for each graph apart from the four well-known exceptions (the Petersen graph, the Coxeter graph and their truncations).'' \cite[\S2]{PSV13}
\end{quote}

\begin{corollary}\label{cor:census}
Every connected cubic Cayley graph on at most $1280$ vertices is $3$-edge-colourable.  More generally, every connected cubic vertex-transitive graph on at most $1280$ vertices other than the Petersen graph and its truncation is $3$-edge-colourable.
\end{corollary}
\begin{proof}
A hamiltonian cubic graph of even order is $3$-edge-colourable.  The Coxeter graph has chromatic index $3$, hence so does its truncation by Lemma~\ref{lem:truncation}; the Petersen graph and its truncation are not $3$-edge-colourable, by Lemma~\ref{lem:truncation} again.  None of the four exceptions is a Cayley graph.
\end{proof}

\subsection{Beyond the census}
Combining Corollary~\ref{cor:census} with Theorem~\ref{thm:NS} gives a concrete list of the smallest groups on which a Cayley snark could conceivably live: a smallest Cayley snark has more than $1280$ vertices, and its group is a non-abelian simple group of order $>1280$, or an almost simple group $T.2$ with $|T|>640$ and $T$ simple, or $S\wr\ZZ_2$ with $|S|>\sqrt{320}$, i.e.\ $|S|\ge60$.  In increasing order the first candidates are $\PGL(2,11)$ (order $1320$), $\PGL(2,13)$ ($2184$), $\PSL(2,17)$ ($2448$), $A_7$ ($2520$), $\PSL(2,19)$ ($3420$), $\PSL(2,16)$ ($4080$), $\PGL(2,17)$ ($4896$), $S_7$ ($5040$), $\PSL(3,3)$ ($5616$), $\operatorname{PSU}(3,3)$ ($6048$), $\PSL(2,23)$ ($6072$), $\PGL(2,19)$ ($6840$), $A_5\wr\ZZ_2$ ($7200$), $\PSL(2,25)$ ($7800$), $M_{11}$ ($7920$), $\PSigmaL(2,16)$ ($8160$), \dots

We wrote a program (\texttt{code/cayley\_snark\_check2.py} in the repository accompanying this survey; an earlier version, \texttt{cayley\_snark\_check.py}, is described below) which, for a given permutation group $G$, enumerates all pairs $(a,x)$ with $a$ an involution, $x$ of odd order and $\langle a,x\rangle=G$, up to conjugation in $G$ and inversion of $x$, and decides whether $\X{G}{a}{x}$ is $3$-edge-colourable.  By Proposition~\ref{prop:trivial} this decides Conjecture~\ref{conj:AZ} for \emph{all} cubic Cayley graphs on $G$.  The decision procedure is Lemma~\ref{lem:pregraph}: for a subgroup $H\le G$ containing no conjugate of $x$, the quotient pregraph $H\backslash X$ has $|G:H|$ vertices, and a $3$-edge-colouring of it---found by the SAT solver CaDiCaL \cite{CaDiCaL} through the python-sat interface, with one Boolean variable ``$e$ has colour~$2$'' and one variable ``$e$ has colour~$1$'' per edge or semi-edge $e$ and the obvious clauses at each vertex---lifts to an $H$-invariant $3$-edge-colouring of $X$.  The lifted colouring is then verified independently on $X$: for every $g\in G$ the program checks that the three edges at $g$ have three different colours and that the colour of $\{g,ga\}$ is the same when computed from $g$ and from $ga$.  Candidate subgroups $H$ are generated once per group---point stabilisers, pointwise and setwise stabilisers of pairs and of small sets, stabilisers of partitions into blocks of equal size, centralisers and normalisers of elements of odd order and of involutions, and abelian subgroups of odd order---(conjugate subgroups give isomorphic quotients, so at most two candidates with the same order and the same multiset of element orders are kept), and for each pair $(a,x)$ those not meeting the conjugacy class of $x$ are tried in decreasing order of size, the full Cayley graph (with a time limit) serving as a last resort; the last resort was never needed.  Generation is tested through the socle $T$ of $G$ ($T=G$ or $|G:T|=2$): $\langle a,x\rangle=G$ if and only if $a\notin T$ and $\langle x,x^a\rangle=T$ (or $\langle a,x\rangle=T$), and a subgroup of $T$ generated by two given elements is all of $T$ as soon as its order exceeds $|T|/m$, where $m$ is the minimal index of a proper subgroup of $T$ (its minimal degree, taken from the Atlas).  Groups are given by permutation generators (constructed from matrices over finite fields, or copied from the Atlas of Finite Group Representations \cite{ATLAS}), and the program checks every group order against the known value; the Petersen graph is rejected as a control, and the SAT encoding for pregraphs was cross-checked against exhaustive search on a thousand random small pregraphs.

The effect of Lemma~\ref{lem:pregraph} is dramatic.  The earlier program, which only used abelian subgroups $H$ of odd order (no semi-edges), reduced $\X{S_8}{a}{x}$ with $x$ of order $15$ to graphs on $4480$ vertices on which the solver timed out, leaving $19$ pairs on $S_8$ undecided; with $H=S_4\wr S_2$ (the stabiliser of a partition into two blocks of four, which contains no element of cycle type $5\cdot3$) the quotient has $35$ vertices and is coloured instantly.  The largest quotient that had to be solved for any group in Table~\ref{tab:computation} has $9\,240$ vertices, whereas the groups have order up to $300\,696$.  Table~\ref{tab:computation} lists, for each group, the number of generating pairs, the orders of $x$ that occur, and the number $\max|\bar V|$ of vertices of the largest quotient pregraph that was needed; all $25\,954$ Cayley graphs were found to be $3$-edge-colourable.

\begin{table}[ht]
\centering
\small
\begin{tabular}{@{}lrrrrrl@{}}
\toprule
group $G$ & $|G|$ & inv.\ classes & pairs $(a,x)$ & orders of $x$ & $\max|\bar V|$ & result\\
\midrule
$A_5$ & 60 & 1 & 6 & 3, 5 & 15 & all colourable\\
$S_5$ & 120 & 2 & 2 & 5 & 5 & all colourable\\
$\mathrm{PSL}(2,7)$ & 168 & 1 & 4 & 3, 7 & 21 & all colourable\\
$\mathrm{PGL}(2,7)$ & 336 & 2 & 5 & 3, 7 & 42 & all colourable\\
$A_6$ & 360 & 1 & 4 & 5 & 10 & all colourable\\
$\mathrm{PSL}(2,8)$ & 504 & 1 & 42 & 3, 7, 9 & 63 & all colourable\\
$\mathrm{PSL}(2,11)$ & 660 & 1 & 14 & 3, 5, 11 & 66 & all colourable\\
$M_{10}$ & 720 & 1 & 0 & -- & -- & no such pairs\\
$\mathrm{PGL}(2,9)$ & 720 & 2 & 10 & 3, 5 & 45 & all colourable\\
$S_6$ & 720 & 3 & 2 & 5 & 10 & all colourable\\
$\mathrm{PSL}(2,13)$ & 1092 & 1 & 45 & 3, 7, 13 & 84 & all colourable\\
$\mathrm{PGL}(2,11)$ & 1320 & 2 & 17 & 3, 5, 11 & 120 & all colourable\\
$\mathrm{PGL}(2,13)$ & 2184 & 2 & 26 & 3, 7, 13 & 168 & all colourable\\
$\mathrm{PSL}(2,17)$ & 2448 & 1 & 60 & 3, 9, 17 & 153 & all colourable\\
$A_7$ & 2520 & 1 & 11 & 5, 7 & 35 & all colourable\\
$\mathrm{PSL}(2,19)$ & 3420 & 1 & 76 & 3, 5, 9, 19 & 180 & all colourable\\
$\mathrm{PSL}(2,16)$ & 4080 & 1 & 204 & 3, 5, 15, 17 & 255 & all colourable\\
$\mathrm{PGL}(2,17)$ & 4896 & 2 & 36 & 3, 9, 17 & 288 & all colourable\\
$S_7$ & 5040 & 3 & 19 & 3, 5, 7 & 42 & all colourable\\
$\mathrm{PSL}(3,3)$ & 5616 & 1 & 20 & 3, 13 & 78 & all colourable\\
$\mathrm{PSU}(3,3)$ & 6048 & 1 & 6 & 7 & 28 & all colourable\\
$\mathrm{PSL}(2,23)$ & 6072 & 1 & 107 & 3, 11, 23 & 276 & all colourable\\
$\mathrm{PGL}(2,19)$ & 6840 & 2 & 58 & 3, 5, 9, 19 & 360 & all colourable\\
$A_5\wr\mathbb{Z}_2$ & 7200 & 3 & 8 & 15 & 25 & all colourable\\
$\mathrm{PSL}(2,25)$ & 7800 & 1 & 134 & 3, 5, 13 & 300 & all colourable\\
$M_{11}$ & 7920 & 1 & 14 & 5, 11 & 55 & all colourable\\
$\mathrm{P\Sigma L}(2,16)$ & 8160 & 2 & 20 & 15, 17 & 68 & all colourable\\
$\mathrm{PSL}(2,27)$ & 9828 & 1 & 186 & 3, 7, 13 & 378 & all colourable\\
$\mathrm{PSL}(3,3).2$ & 11232 & 2 & 31 & 3, 13 & 156 & all colourable\\
$G_2(2)=\mathrm{PSU}(3,3).2$ & 12096 & 2 & 21 & 3, 7 & 378 & all colourable\\
$\mathrm{PGL}(2,23)$ & 12144 & 2 & 72 & 3, 11, 23 & 528 & all colourable\\
$\mathrm{PSL}(2,29)$ & 12180 & 1 & 219 & 3, 5, 7, 15, 29 & 435 & all colourable\\
$\mathrm{PSL}(2,31)$ & 14880 & 1 & 160 & 3, 5, 15, 31 & 480 & all colourable\\
$\mathrm{PGL}(2,25)$ & 15600 & 2 & 90 & 3, 5, 13 & 300 & all colourable\\
$\mathrm{PSL}(2,25).2_3$ & 15600 & 1 & 0 & -- & -- & no such pairs\\
$\mathrm{P\Sigma L}(2,25)$ & 15600 & 3 & 30 & 13 & 26 & all colourable\\
$\mathrm{PGL}(2,27)$ & 19656 & 2 & 123 & 3, 7, 13 & 378 & all colourable\\
$A_8$ & 20160 & 2 & 33 & 5, 7, 15 & 35 & all colourable\\
$\mathrm{PSL}(3,4)$ & 20160 & 1 & 60 & 5, 7 & 105 & all colourable\\
$\mathrm{PGL}(2,29)$ & 24360 & 2 & 146 & 3, 5, 7, 15, 29 & 840 & all colourable\\
$\mathrm{PSL}(2,37)$ & 25308 & 1 & 378 & 3, 9, 19, 37 & 684 & all colourable\\
$\mathrm{PSU}(4,2)=\mathrm{PSp}(4,3)$ & 25920 & 2 & 42 & 5, 9 & 120 & all colourable\\
$Sz(8)$ & 29120 & 1 & 210 & 5, 7, 13 & 455 & all colourable\\
$\mathrm{PGL}(2,31)$ & 29760 & 2 & 113 & 3, 5, 15, 31 & 960 & all colourable\\
$\mathrm{PSL}(2,32)$ & 32736 & 1 & 930 & 3, 11, 31, 33 & 1023 & all colourable\\
$\mathrm{PSL}(2,41)$ & 34440 & 1 & 380 & 3, 5, 7, 21, 41 & 861 & all colourable\\
$\mathrm{PSL}(2,43)$ & 39732 & 1 & 490 & 3, 7, 11, 21, 43 & 924 & all colourable\\
$\mathrm{PSL}(3,4).2$ (field) & 40320 & 2 & 11 & 5, 7 & 105 & all colourable\\
$\mathrm{PSL}(3,4).2$ (graph) & 40320 & 2 & 68 & 3, 5, 7 & 315 & all colourable\\
$\mathrm{PSL}(3,4).2$ (graph-field) & 40320 & 2 & 45 & 5, 7 & 105 & all colourable\\
$S_8$ & 40320 & 4 & 47 & 5, 7, 15 & 35 & all colourable\\
$\mathrm{PGL}(2,37)$ & 50616 & 2 & 242 & 3, 9, 19, 37 & 1368 & all colourable\\
$\mathrm{PSU}(4,2).2=W(E_6)$ & 51840 & 4 & 52 & 5, 9 & 40 & all colourable\\
$\mathrm{PSL}(2,47)$ & 51888 & 1 & 437 & 3, 23, 47 & 1128 & all colourable\\
$\mathrm{PSL}(2,7)\wr\mathbb{Z}_2$ & 56448 & 3 & 6 & 21 & 42 & all colourable\\
$\mathrm{PSL}(2,49)$ & 58800 & 1 & 480 & 3, 5, 7, 25 & 1200 & all colourable\\
$\mathrm{PSU}(3,4)$ & 62400 & 1 & 64 & 3, 5, 13, 15 & 195 & all colourable\\
$\mathrm{PGL}(2,41)$ & 68880 & 2 & 250 & 3, 5, 7, 21, 41 & 1680 & all colourable\\
$\mathrm{PSL}(2,53)$ & 74412 & 1 & 780 & 3, 9, 13, 27, 53 & 1431 & all colourable\\
$\mathrm{PGL}(2,43)$ & 79464 & 2 & 325 & 3, 7, 11, 21, 43 & 1848 & all colourable\\
$M_{12}$ & 95040 & 2 & 66 & 3, 5, 11 & 396 & all colourable\\
$\mathrm{PSL}(2,59)$ & 102660 & 1 & 926 & 3, 5, 15, 29, 59 & 1770 & all colourable\\
$\mathrm{PGL}(2,47)$ & 103776 & 2 & 288 & 3, 23, 47 & 2208 & all colourable\\
$\mathrm{PSL}(2,61)$ & 113460 & 1 & 1023 & 3, 5, 15, 31, 61 & 1860 & all colourable\\
$\mathrm{PGL}(2,49)$ & 117600 & 2 & 324 & 3, 5, 7, 25 & 2400 & all colourable\\
$\mathrm{PSL}(2,49).2_3$ & 117600 & 1 & 0 & -- & -- & no such pairs\\
$\mathrm{P\Sigma L}(2,49)$ & 117600 & 3 & 80 & 5, 25 & 50 & all colourable\\
$\mathrm{PSU}(3,4).2$ & 124800 & 2 & 200 & 3, 5, 13, 15 & 1248 & all colourable\\
$\mathrm{PSU}(3,5)$ & 126000 & 1 & 51 & 5, 7 & 525 & all colourable\\
$\mathrm{PGL}(2,53)$ & 148824 & 2 & 506 & 3, 9, 13, 27, 53 & 2808 & all colourable\\
$\mathrm{PSL}(2,67)$ & 150348 & 1 & 1216 & 3, 11, 17, 33, 67 & 2244 & all colourable\\
$J_1$ & 175560 & 1 & 544 & 3, 5, 7, 11, 15, 19 & 9240 & all colourable\\
$\mathrm{PSL}(2,71)$ & 178920 & 1 & 1124 & 3, 5, 7, 9, 35, 71 & 2556 & all colourable\\
$A_9$ & 181440 & 2 & 257 & 3, 5, 7, 9, 15 & 280 & all colourable\\
$M_{12}.2$ & 190080 & 3 & 55 & 3, 5, 11 & 792 & all colourable\\
$\mathrm{PSL}(2,73)$ & 194472 & 1 & 1248 & 3, 9, 37, 73 & 2664 & all colourable\\
$\mathrm{PGL}(2,59)$ & 205320 & 2 & 623 & 3, 5, 15, 29, 59 & 3480 & all colourable\\
$\mathrm{PGL}(2,61)$ & 226920 & 2 & 674 & 3, 5, 15, 31, 61 & 3720 & all colourable\\
$\mathrm{PSL}(2,79)$ & 246480 & 1 & 1258 & 3, 5, 13, 39, 79 & 3120 & all colourable\\
$\mathrm{PSU}(3,5).2$ & 252000 & 2 & 102 & 3, 5, 7 & 525 & all colourable\\
$A_6\wr\mathbb{Z}_2$ & 259200 & 3 & 20 & 5, 15 & 90 & all colourable\\
$\mathrm{PSL}(2,64)$ & 262080 & 1 & 3858 & 3, 5, 7, 9, 13, 21, 63, 65 & 4095 & all colourable\\
\bottomrule
\end{tabular}

\caption{Computer verification of Conjecture~\ref{conj:AZ}.  For each group the number of generating pairs $(a,x)$ with $a^2=1$, $\ord(x)$ odd, up to conjugation and inversion of $x$, is listed, together with the orders of $x$ that occur and the number of vertices of the largest quotient pregraph that had to be solved; every corresponding cubic Cayley graph was found to be $3$-edge-colourable.  $\PSL(3,4).2$ (field), (graph), (graph-field) are the three index-two extensions of $\PSL(3,4)$, by a field, a graph and a graph-field automorphism respectively.  Groups of order at most $1280$ are included as a consistency check against \cite{PSV13}.}
\label{tab:computation}
\end{table}

\begin{theorem}\label{thm:computation}
Conjecture~\ref{conj:AZ} holds for every cubic Cayley graph on each of the 85 groups listed in Table~\ref{tab:computation}; in particular for every non-abelian simple group of order at most $285\,852$, for $\PSL(2,q)$ with $q\le 83$, for $\PGL(2,q)$ with $q\le 67$, for $A_n$ with $n\le 9$ and $S_n$ with $n\le 8$, and for the groups $M_{10}$, $\mathrm{PSL}(3,3)$, $\mathrm{PSU}(3,3)$, $A_5\wr\mathbb{Z}_2$, $M_{11}$, $\mathrm{P\Sigma L}(2,16)$, $\mathrm{PSL}(3,3).2$, $G_2(2)=\mathrm{PSU}(3,3).2$, $\mathrm{P\Sigma L}(2,25)$, $\mathrm{PSL}(2,25).2_3$, $\mathrm{PSL}(3,4)$, $\mathrm{PSU}(4,2)=\mathrm{PSp}(4,3)$, $Sz(8)$, $\mathrm{PSL}(3,4).2$ (field), $\mathrm{PSL}(3,4).2$ (graph), $\mathrm{PSL}(3,4).2$ (graph-field), $\mathrm{PSU}(4,2).2=W(E_6)$, $\mathrm{PSL}(2,7)\wr\mathbb{Z}_2$, $\mathrm{PSU}(3,4)$, $M_{12}$, $\mathrm{PSL}(2,49).2_3$, $\mathrm{P\Sigma L}(2,49)$, $\mathrm{PSU}(3,4).2$, $\mathrm{PSU}(3,5)$, $J_1$, $M_{12}.2$, $\mathrm{PSU}(3,5).2$, $A_6\wr\mathbb{Z}_2$.
\end{theorem}

Two of the groups in the table, $M_{10}$ and $\PSL(2,25).2_3$, have all their involutions inside the socle and are therefore not generated by an involution and an element of odd order at all; the only cubic Cayley graphs on them are the trivially colourable ones of Proposition~\ref{prop:trivial}.  The groups permitted by Theorem~\ref{thm:NS} for a smallest Cayley snark are the non-abelian simple groups, their extensions $T.2$ of index two and the wreath products $S\wr\ZZ_2$.  Listing these by order (from the orders of the simple groups; note that $\PSL(2,8)$, $Sz(8)$ and $\PSL(2,32)$ have no extension of index two, while $\PSL(2,25)$, $\PSL(2,49)$ and $\PSL(3,4)$ have three each), one finds that every such group of order less than $352\,440$ appears in Table~\ref{tab:computation}, the first ones missing being $\mathrm{PSL}(2,89)$ (order $352\,440$) and $S_9$ (order $362\,880$).  Combining Theorem~\ref{thm:NS} with Corollary~\ref{cor:census} and Table~\ref{tab:computation} therefore gives:

\begin{corollary}\label{cor:bound}
A Cayley snark of smallest order, if one exists, has more than $352\,439$ vertices.  Its group is a non-abelian simple group, or a group $T.2$ or $S\wr\ZZ_2$ as in Theorem~\ref{thm:NS}\textup{(ii)}, of order at least $352\,440$ and not listed in Table~\ref{tab:computation}.
\end{corollary}

We stress that these computations, like the census, are evidence rather than insight: the even $2$-factors found by the solver do not seem to follow any pattern that would suggest a general construction.  One structural remark can be made, though.  In every case a colouring invariant under a fairly large subgroup was found---the quotient never had more than $9\,240$ vertices---and for the groups $\PSL(2,q)$ with $x$ in a non-split torus the $(q+1)$-vertex quotient by a Borel subgroup sufficed for the great majority of pairs once $q\ge 25$ (Example~\ref{ex:projline}); but for small $q$, and for $q$ even, it does not, so ``the'' natural quotient is not the answer.  Another instructive family is $\PGL(2,q)$ with $\ord(x)=3$ and $\ord(ax)=q+1$: the quotients by the dihedral subgroups of order $2(q+1)$ (admissible when $3\nmid q+1$) were never $3$-edge-colourable---for $q=61$ the solver proved all $30$ such quotients, on $1830$ and $3660$ vertices, uncolourable---whereas the quotient by the unipotent subgroup of order $q$, or by a Borel subgroup when $3\nmid q-1$, always was.

\section{Related problems on flows in Cayley graphs}\label{sec:related}

The flow formulation of Conjecture~\ref{conj:AZ} has stimulated work on the stronger $3$-flow property.  By Mader's theorem (Theorem~\ref{thm:mader}), Tutte's $3$-flow conjecture restricted to vertex-transitive graphs asserts that every vertex-transitive graph of valency at least $4$ admits a nowhere-zero $3$-flow.  This is known for Cayley graphs on abelian groups (Poto\v{c}nik, \v{S}koviera and \v{S}krekovski \cite{PSS05}), for Cayley graphs of valency at least $4$ on groups with a certain Sylow $2$-subgroup structure (N\'an\'asiov\'a and \v{S}koviera \cite{NanasiovaSkoviera09}), for graphs admitting a solvable arc-transitive group of automorphisms (Li and Zhou \cite{LiZhou16}), for Cayley graphs on supersolvable groups (Zhang and Zhou \cite{ZhangZhou24}), for nilpotently vertex-transitive graphs \cite{ZhangZhou22b}, for various dihedral and generalised dihedral and quaternion groups \cite{YangLi11,AhanjidehIranmanesh19}, and for Cayley graphs on solvable groups of twice square-free order \cite{AhanjidehKovacs26}.  The general case for solvable groups, which would be the $3$-flow analogue of Theorem~\ref{thm:ALZ}, is open.  These results do not bear on the cubic case (a cubic graph never has a nowhere-zero $3$-flow), but their proofs use the same quotient-and-lift technique of \S\ref{sec:cover}.

\subsection{The cycle-double-cover proof and an eight-point palette}\label{sec:cdc-palette}

In July 2026 OpenAI announced a proof of the cycle double cover conjecture \cite{OpenAICDC26}; expositions by Geelen \cite{Geelen26} and Oum \cite{Oum26} give independent presentations of the argument.  Its starting point for a loopless bridgeless cubic graph $X$ is a nowhere-zero flow $f\colon E(X)\to\Gamma$, where $\Gamma=\mathbb F_2^3$.  It constructs a two-element set $P_e\subset\Gamma$ for every edge such that
\begin{equation}\label{eq:cdc-parity}
 m_v(s):=|\{e\ni v:s\in P_e\}|\in\{0,2\}
 \qquad(v\in V(X),\ s\in\Gamma).
\end{equation}
The even subgraphs $M_s=\{e:s\in P_e\}$ then form a cycle double cover.  More explicitly, after ordering the three edges at $v$ as $a,b,c$ and writing $f(a)=x$, $f(b)=y$, $f(c)=x+y$, the proof chooses local offsets $g_{v,a}=0$, $g_{v,b}=x$, $g_{v,c}=0$.  It proves by a duality argument that the affine system
\begin{equation}\label{eq:cdc-system}
 t_u+t_v+\epsilon_e f(e)=g_{u,e}+g_{v,e}
 \qquad(e=uv)
\end{equation}
has a solution with $t_v\in\Gamma$ and $\epsilon_e\in\mathbb F_2$.  The sets
\[
 P_e=\{t_v+g_{v,e},\ t_v+g_{v,e}+f(e)\}
\]
are consequently independent of the chosen endpoint and satisfy \eqref{eq:cdc-parity}.

There is a precise way to measure how far this relaxed colouring is from a Tait colouring.  For a family $P=(P_e)$ satisfying \eqref{eq:cdc-parity}, define its \emph{palette graph} $H(P)$ to be the simple graph with vertex set $\Gamma$ and edge set $\{P_e:e\in E(X)\}$ (repetitions are suppressed).

\begin{proposition}[palette extraction]\label{prop:palette}
Let $X$ be a loopless cubic graph and let $P_e\subset\mathbb F_2^3$ be two-element sets satisfying \eqref{eq:cdc-parity}.  If $H(P)$ is $4$-vertex-colourable, then $X$ is $3$-edge-colourable.  Conversely, if $X$ is $3$-edge-colourable, one can start the cycle-double-cover construction with a flow $f$ for which every resulting palette graph is $4$-colourable.  Hence $X$ is $3$-edge-colourable if and only if some output of the construction has a $4$-colourable palette graph.
\end{proposition}
\begin{proof}
Identify the four vertex colours with the elements of $\mathbb F_2^2$, and let $c\colon\Gamma\to\mathbb F_2^2$ be a proper colouring of $H(P)$.  For $P_e=\{p_e,q_e\}$ put
\[
 h(e)=c(p_e)+c(q_e).
\]
This is nonzero because $c(p_e)\ne c(q_e)$.  At every vertex $v$,
\[
 \sum_{e\ni v}h(e)
 =\sum_{s\in\Gamma}m_v(s)c(s)=0,
\]
since every $m_v(s)$ is zero or two.  Thus $h$ is a nowhere-zero $\mathbb F_2^2$-flow, whose three nonzero values give a $3$-edge-colouring of the cubic graph.

Conversely, encode the three classes of a $3$-edge-colouring by the three nonzero elements of a two-dimensional subspace $W<\Gamma$.  This gives a nowhere-zero $\Gamma$-flow $f$ with image in $W$.  Choose a linear map $\pi\colon\Gamma\to\mathbb F_2^2$ whose restriction to $W$ is injective.  In every output of the construction the two ends of the palette edge $P_e$ differ by $f(e)\in W\setminus\{0\}$, so $\pi$ is a proper $4$-colouring of $H(P)$.
\end{proof}

One particularly concrete sufficient target follows.  If the complement of $H(P)$ has a perfect matching, give the two vertices in each matched pair the same one of four colours; then Proposition~\ref{prop:palette} gives a Tait colouring of $X$.  Thus a possible route to Conjecture~\ref{conj:AZ} is to use the Cayley structure to choose the initial $\mathbb F_2^3$-flow and the affine solution so that the palette is $4$-colourable---or, more strongly, so that its complement has a perfect matching.

The converse half of Proposition~\ref{prop:palette} also exposes the danger of circularity: a rank-two starting flow is already a Tait colouring in disguise.  The construction itself guarantees the parity condition \eqref{eq:cdc-parity}, not the global chromatic bound on $H(P)$.  Nor can one simply choose the starting flow to inherit the full Cayley symmetry.

\begin{lemma}\label{lem:no-invariant-eight-flow}
Let $X=\X{G}{a}{x}$ with $\ord(x)>2$.  No nowhere-zero $\mathbb F_2^k$-flow on $X$ is invariant under the left regular action of $G$, for any $k\ge1$.
\end{lemma}
\begin{proof}
An invariant assignment is constant, say $\alpha$, on the orbit of $a$-edges and constant, say $\xi$, on the orbit of $x$-edges.  At each vertex there is one edge of the first kind and two of the second.  In characteristic two the flow equation is therefore $\alpha+\xi+\xi=\alpha=0$, contrary to nowhere-zeroness.
\end{proof}

Full rank by itself does not remove the circularity.  In fact the exact reformulation in Proposition~\ref{prop:palette} remains exact after full rank is imposed.

\begin{proposition}[full-rank tetrahedral palette]\label{prop:full-rank-palette}
Let $X$ be a connected loopless cubic graph on $n\ge6$ vertices.  Then $X$ is $3$-edge-colourable if and only if the cycle-double-cover construction has an output whose starting $\mathbb F_2^3$-flow has rank three and whose palette pairs are supported on exactly four affinely independent elements of $\mathbb F_2^3$.
\end{proposition}
\begin{proof}
The reverse implication is Proposition~\ref{prop:palette}.  For the forward implication, let $h\colon E(X)\to\mathbb F_2^2\setminus\{0\}$ be the flow encoded by a $3$-edge-colouring and put $m=|E(X)|=3n/2$.  Consider the homogeneous linear system
\[
 r_u+r_v+\delta_e h(e)=0\qquad(e=uv),
\]
with $r_v\in\mathbb F_2^2$ and $\delta_e\in\mathbb F_2$.  It has $2n+m$ scalar variables and at most $2m$ independent equations, so its solution space has dimension at least $2n-m=n/2\ge3$.  The constant solutions $r_v=r$, $\delta_e=0$ form only a two-dimensional subspace; choose a nonconstant solution.

Index four affinely independent points $q_z\in\mathbb F_2^3$ by $z\in\mathbb F_2^2$.  For $e=uv$, the four indices split into two pairs $\{z,z+h(e)\}$.  Let $P_e$ be the pair of corresponding $q$'s which does not contain $q_{r_u}$.  The equation says $r_v\in\{r_u,r_u+h(e)\}$, so the same pair avoids $q_{r_v}$ and $P_e$ is well-defined.  At $v$ the three incident values of $h$ are the three nonzero elements of $\mathbb F_2^2$; hence the three $P_e$ are precisely the edges of the triangle on the three points other than $q_{r_v}$.  These local triangles have exactly the form used in the cycle-double-cover construction, and their agreement across edges gives a solution of its compatibility system.  Put $f(e)=p_e+q_e$ when $P_e=\{p_e,q_e\}$.  The triangle condition makes $f$ a nowhere-zero $\mathbb F_2^3$-flow.  Since $r$ is nonconstant, at least two different local triangles occur; their union is connected and uses all four affinely independent points, so their edge differences span $\mathbb F_2^3$.  Thus $f$ has rank three.
\end{proof}

For a fixed starting flow, the freedom in the paper's affine system has a simple intrinsic description.  Define the \emph{weighted cut code}
\begin{equation}\label{eq:weighted-cut-code}
 \mathcal K_f=\left\{\delta\in\mathbb F_2^{E(X)}:
 \sum_{e\in C}\delta_e f(e)=0\text{ for every cycle }C\text{ of }X\right\}.
\end{equation}
Equivalently, for each of the three coordinates $f_j$ of $f$, the edge set with indicator $\delta_e f_j(e)$ must be a cut.  This follows from the orthogonality of the binary cycle and cut spaces, and explains the terminology.

\begin{proposition}[compatibility-space rigidity]\label{prop:cdc-rigidity}
Fix a nowhere-zero $\mathbb F_2^3$-flow $f$ on a connected loopless cubic graph and fix the local data $g_{v,e}$ in the cycle-double-cover construction.  The nullity of the homogeneous compatibility system associated with \eqref{eq:cdc-system} is
\[
 3+\dim\mathcal K_f.
\]
Modulo a common translation of every palette point by an element of $\mathbb F_2^3$, its solution space is therefore parameterised by $\mathcal K_f$.  In particular, if $\mathcal K_f=0$, all outputs for these fixed local data have isomorphic palette graphs.
\end{proposition}
\begin{proof}
Subtract two solutions of \eqref{eq:cdc-system}.  Writing their differences as $s_v\in\mathbb F_2^3$ and $\delta_e\in\mathbb F_2$ gives
\begin{equation}\label{eq:cdc-homogeneous}
 s_u+s_v=\delta_e f(e)\qquad(e=uv).
\end{equation}
The right-hand side is the coboundary of a vertex labelling $s$ if and only if its sum around every cycle is zero, which is exactly the condition $\delta\in\mathcal K_f$.  Indeed, necessity follows by summing \eqref{eq:cdc-homogeneous} around a cycle; conversely, fix $s$ at one vertex and define it by path sums, which are path-independent by the cycle condition.  For each $\delta\in\mathcal K_f$ the resulting $s$ is unique up to one common element of $\mathbb F_2^3$, accounting for the other three dimensions.  Such a common translation translates every palette pair and hence preserves the palette graph up to isomorphism.
\end{proof}

Thus a fixed flow with $\mathcal K_f=0$ offers no effective affine choice at all: one must change the flow or the local triangles.  More generally, any extra matching equations can only be met through the evaluation of the weighted cut code (together with the irrelevant global translations).  A large code is not sufficient, but the proposition isolates exactly where the cycle-double-cover construction can supply non-circular flexibility.

There is nevertheless a more structured way in which the Cayley generators can enter the construction.  Let $Q$ be the multigraph obtained from $X=\X{G}{a}{x}$ by contracting every $x$-cycle to a vertex; its edges are the $a$-edges of $X$.  Thus $Q$ is a connected, vertex-transitive, $s$-regular loopless multigraph, where $s=\ord(x)$.

\begin{proposition}[generator-separated $8$-flow]\label{prop:separated-eight-flow}
Suppose that $s$ is odd and that $Q$ has a nowhere-zero $\mathbb F_2^2$-flow.  Then $X$ has a rank-three nowhere-zero flow $f\colon E(X)\to\mathbb F_2\oplus\mathbb F_2^2$ such that every $a$-edge has value $(0,p)$ with $p\ne0$, while every $x$-edge has value $(1,y)$.
\end{proposition}
\begin{proof}
Let $\phi\colon E(Q)\to\mathbb F_2^2\setminus\{0\}$ be such a flow.  Around an $x$-cycle write the values of its incident $a$-edges, in cyclic order, as $p_0,\ldots,p_{s-1}$.  The flow equation in $Q$ says $\sum_i p_i=0$, so there are $y_i\in\mathbb F_2^2$ satisfying
\[
 y_{i-1}+y_i=p_i\qquad(i\bmod s):
\]
choose one $y_{s-1}$ and continue recursively; the last equation is exactly $\sum_i p_i=0$.  Assign $(0,p_i)$ to the corresponding $a$-edge and $(1,y_i)$ to the $i$th $x$-edge.  The displayed equation is precisely the flow equation at the corresponding vertex of $X$, and every assigned value is nonzero.

It remains to check rank.  If the three nonzero elements of $\mathbb F_2^2$ occur $n_1,n_2,n_3$ times at a vertex of $Q$, the flow equation forces $n_1,n_2,n_3$ to have the same parity.  Their sum is the odd number $s$, so all three are odd and in particular positive.  The $a$-edge values therefore span $\{0\}\oplus\mathbb F_2^2$, and any $x$-edge value supplies the missing first coordinate.
\end{proof}

The hypothesis of Proposition~\ref{prop:separated-eight-flow} is automatic when $s\ge5$ and $Q$ is simple: Mader's theorem makes $Q$ $s$-edge-connected, and Jaeger's theorem and Tutte's group-flow theorem give the required $\mathbb F_2^2$-flow.  Writing $H=\langle x\rangle$, the neighbours of the vertex $H$ are the cosets $haH$ with $h\in H$, so two coincide exactly when their quotient lies in $H\cap aHa$.  Thus $Q$ is simple exactly when $H\cap aHa=1$.  In particular this holds when $s$ is prime and $G$ is non-solvable, since otherwise $aHa=H$, so $a$ normalises $H$ and $G=\langle a,x\rangle$ is solvable.

This separates the palette into two four-point layers.  In any output $P$, an $a$-edge is a pair within one layer and an $x$-edge is a pair crossing the layers.  Let $B_x(P)$ be the bipartite graph on the two layers formed by the palette pairs belonging to $x$-edges.  If the bipartite complement of $B_x(P)$ has a perfect matching, colour the two endpoints of each matched pair alike.  All within-layer pairs and all unmatched crossing pairs are then properly coloured, so Proposition~\ref{prop:palette} gives a Tait colouring.  The paper's method has therefore reduced this structured case to a four-by-four Hall problem.  What it does not yet supply is a choice of $\phi$, of the lifts $y_i$, and of the compatibility-system solution which forces Hall's condition.

The twenty-four possible matchings admit a useful exact classification.  Put $A=\mathbb F_2^2$, so that the layers are $\{0\}\times A$ and $\{1\}\times A$.

\begin{proposition}[linearised matching target]\label{prop:linearised-matching}
Every perfect matching between the two layers has the form
\[
 \mathcal M_{L,c}=\bigl\{\{(0,z),(1,Lz+c)\}:z\in A\bigr\}
\]
for unique $L\in\operatorname{GL}(2,2)$ and $c\in A$.  Let an $x$-edge have flow value $(1,y_e)$ and write its palette pair as $\{(0,z_e),(1,z_e+y_e)\}$.  This pair belongs to $\mathcal M_{L,c}$ exactly when
\begin{equation}\label{eq:matching-collision}
 (I+L)z_e=y_e+c.
\end{equation}
For each of the twelve choices with $L$ of order two, requiring that no $x$-edge pair belong to $\mathcal M_{L,c}$ amounts to adding ordinary affine linear equations to the compatibility system~\eqref{eq:cdc-system}.
\end{proposition}
\begin{proof}
The affine group on $A$ has order $4\cdot6=24$ and induces all $24$ permutations of its four points, proving the first statement.  The equality of the upper endpoints of the two displayed pairs gives \eqref{eq:matching-collision}.  If $L$ has order two, $N=I+L$ has rank one and $K:=\ker N=\operatorname{im}N$.  Let $\rho\colon A\to\mathbb F_2$ be the nonzero functional with kernel $K$, and let $\lambda\colon K\to\mathbb F_2$ be its nonzero functional.  Then $\lambda(Nz)=\rho(z)$.  When $y_e+c\notin K$, collision is impossible.  Otherwise \eqref{eq:matching-collision} is equivalent to
\[
 \rho(z_e)=\lambda(y_e+c),
\]
so avoiding collision is the single affine equation $\rho(z_e)=1+\lambda(y_e+c)$.

It remains only to see that this really is linear in the paper's variables.  At an endpoint $v$, write $t_v+g_{v,e}=(b,w)\in\mathbb F_2\oplus A$.  Since the other endpoint of the pair is obtained by adding $(1,y_e)$, its layer-zero endpoint has lower coordinate
\[
 z_e=w+b y_e,
\]
an affine linear function of $t_v$.  Thus each required non-collision equation can be appended directly to \eqref{eq:cdc-system}.
\end{proof}

Proposition~\ref{prop:cdc-rigidity} turns consistency of this augmented system into a smaller, exact target.  Let $U_f$ be the homogeneous solution space of the original compatibility system and let $T<U_f$ be its three-dimensional subspace of global translations.

\begin{proposition}[transvection syndrome criterion]\label{prop:transvection-syndrome}
Fix $f$, the local data, and a transvection matching $\mathcal M_{L,c}$.  Let $D$ be the set of $x$-edges for which $y_e+c\in\operatorname{im}(I+L)$; these are exactly the edges on which collision is possible.  Write the corresponding non-collision equations as
\[
 A\mathbf z=b\quad\text{in }\mathbb F_2^D,
\]
where $\mathbf z=(t,\epsilon)$ denotes the variables of the compatibility system, and fix any solution $\mathbf z_0$ of that original system.  Then the augmented system is consistent if and only if
\begin{equation}\label{eq:syndrome-criterion}
 [b+A\mathbf z_0]\in
 \operatorname{im}\left(
  \overline A:\mathcal K_f\cong U_f/T
  \longrightarrow \mathbb F_2^D/A(T)
 \right).
\end{equation}
In particular,
\[
 \dim\operatorname{im}\overline A
 =\operatorname{rank}(A|_{U_f})-\operatorname{rank}(A|_T)
 \le \dim\mathcal K_f.
\]
Failure is certified by a linear functional on $\mathbb F_2^D$ which annihilates $A(U_f)$ but is nonzero on $b+A\mathbf z_0$.
\end{proposition}
\begin{proof}
Every solution of the original system is $\mathbf z_0+u$ with $u\in U_f$.  It satisfies the extra equations exactly when
\[
 Au=b+A\mathbf z_0.
\]
Thus consistency is equivalent to $b+A\mathbf z_0\in A(U_f)$.  Quotienting the target by $A(T)$ and using $U_f/T\cong\mathcal K_f$ from Proposition~\ref{prop:cdc-rigidity} gives \eqref{eq:syndrome-criterion} and the rank formula.  The final statement is the standard dual characterisation of membership in a linear subspace.
\end{proof}

The explicit equation in the proof of Proposition~\ref{prop:linearised-matching} shows also that $A$ depends on $f$ and the chosen matching but not on the local data $g_{v,e}$; changing the local triangles changes the syndrome $b+A\mathbf z_0$, not the available image.  Thus flow selection and local-triangle selection have distinct roles.  The last sentence of Proposition~\ref{prop:transvection-syndrome} identifies precisely how the paper's duality methodology would have to be strengthened: its original argument rules out inconsistency certificates for the compatibility equations, whereas one must now rule out certificates carrying an additional functional supported on the dangerous edges.  Merely counting variables cannot do this, since the target syndrome may lie outside the usually much smaller image of $\overline A$.

This splits the Hall strategy into four parallel matchings ($L=I$), twelve transvection matchings (the linear case of Proposition~\ref{prop:linearised-matching}), and eight matchings for which $L$ has order three.  In the last case $I+L$ is invertible, so \eqref{eq:matching-collision} forbids one point rather than one affine line and the condition is no longer linear.  The parallel case can be understood completely, and shows precisely where the cycle-double-cover argument adds no leverage.

\begin{proposition}[parallel-matching criterion]\label{prop:parallel-matching}
Retain the quotient flow $\phi$ of Proposition~\ref{prop:separated-eight-flow}.  At a vertex of $Q$, list its incident values cyclically as $p_0,\ldots,p_{s-1}$ and put
\[
 S=\{0,p_0,p_0+p_1,\ldots,p_0+\cdots+p_{s-2}\}\subseteq A.
\]
For any fixed $c\in A$, the lifts can be chosen so that the parallel matching $\mathcal M_{I,c}$ misses every palette pair if and only if $S\ne A$ at every vertex of $Q$.  Whenever this happens, the Tait flow is already given directly by
\[
 h(e)=p_e\quad(e\text{ an }a\text{-edge}),\qquad
 h(e)=y_e+c\quad(e\text{ an }x\text{-edge}).
\]
If $s=5$, the local condition $S\ne A$ says exactly that the value occurring three times among $p_0,\ldots,p_4$ occurs in three cyclically consecutive positions.
\end{proposition}
\begin{proof}
All solutions of $y_{i-1}+y_i=p_i$ around the cycle are obtained by translating the set of $y_i$ by a common element of $A$, and that set is a translate of $S$.  It can be made to avoid $c$ exactly when $S\ne A$.  By \eqref{eq:matching-collision} with $L=I$, an $x$-edge pair belongs to $\mathcal M_{I,c}$ exactly when $y_e=c$, independently of the compatibility-system solution.  The displayed $h$ is nonzero, and at a vertex its sum is $p_i+(y_{i-1}+c)+(y_i+c)=0$.

For $s=5$, the flow equation makes the multiplicities of the three nonzero elements of $A$ all odd, hence they are $3,1,1$.  Write the repeated value as $r$ and the other two as $u,v$, so $r=u+v$.  If the three $r$'s are consecutive, a direct partial-sum calculation omits one element of $A$.  Otherwise the two exceptional entries separate the three $r$'s into nonempty blocks; up to reversal and cyclic shift the word is $u,r,v,r,r$, whose partial sums visit $u,v,0,r$, all of $A$.
\end{proof}

Thus the four parallel targets are merely direct Tait-flow criteria in different coordinates.  The twelve transvection targets are the part of the paper's methodology that remains potentially useful: they ask whether the duality proof of consistency for \eqref{eq:cdc-system} can survive a particular family of extra linear equations.

We made a small diagnostic computation, reproducible with \nolinkurl{code/cdc_palette_experiment.py}.  On $K_4$, five distinct full-rank flows gave $40$ labelled outputs, all with palette chromatic number $4$.  For three distinct full-rank flows on the Petersen graph, one choice of local orders per flow led to affine solution spaces of dimensions $3,3,4$; all $32$ resulting palettes had chromatic number $5$, as they must if Proposition~\ref{prop:palette} is to pass this negative control.  For the colourable graph $\Cay(S_5,\{(0\,1),(0\,1\,2\,3\,4),(0\,4\,3\,2\,1)\})$, twenty distinct full-rank flows gave $600$ affine solutions (counted as labelled outputs): $240$ had palette chromatic number $7$ and $360$ had palette chromatic number $8$, with none $4$-colourable.  A direct SAT search does find a rank-three flow with a tetrahedral palette on this graph, as Proposition~\ref{prop:full-rank-palette} predicts, but that search simply solves an equivalent form of the Tait-colouring problem.  A second SAT search finds a generator-separated rank-three flow and an output certified by a transvection matching, while forcing all four lower-coordinate differences to occur on $x$-edges; hence none of the four parallel matchings can certify that output.  This confirms that the linearised target is strictly more flexible than Proposition~\ref{prop:parallel-matching}, although the search again uses the desired colouring condition rather than proving it.  The experiment shows that merely taking an arbitrary $8$-flow and solving the paper's linear system need not reveal the desired colouring, even for a small Cayley graph.  The new method therefore supplies sharp reformulations and a potentially useful family of augmented linear systems, but not at present a proof of the Alspach--Zhang conjecture.

For that separated transvection witness the compatibility system has nullity $8$, so $\dim\mathcal K_f=5$.  Of its $120$ $x$-edges, $40$ are dangerous for the displayed matching.  The restriction map $A|_{U_f}$ has rank $5$, its global-translation part has rank $1$, and hence the induced map $\overline A$ in Proposition~\ref{prop:transvection-syndrome} has rank $4$.  Thus the forty collision equations reduce, after compatibility and translation are taken into account, to a four-dimensional effective syndrome which the witness hits.  The outstanding proof problem is to force that hit without first knowing a Tait colouring.

As a check against selection bias, the script also constructs twelve distinct quotient $\mathbb F_2^2$-flows, chooses lifts and local data independently of any colouring target, and tests all twelve transvection matchings after solving the original compatibility system.  The resulting weighted cut-code dimensions were $0$ once, $1$ eight times, and $2$ three times; none of the $144$ augmented systems was consistent.  This small sample is not evidence against the conjecture---the graph is colourable and the preceding witness succeeds---but it shows that the successful dimension-five flow is special and that an arbitrary application of the cycle-double-cover solver is unlikely to find it.

\section{Open problems}\label{sec:open}

\begin{enumerate}[leftmargin=2em]
\item (Conjecture~\ref{conj:AZ} for simple groups.)  Is $\X{G}{a}{x}$ $3$-edge-colourable for every non-abelian simple group $G=\langle a,x\rangle$, $a^2=1$, $\ord(x)$ odd?  Even the case $G=\PSL(2,q)$ or $G=A_n$ for all $q$, $n$ is open, although Corollary~\ref{cor:2s3} covers those generating pairs with $\ord(ax)=3$.
\item (The bi-Cayley case.)  For a non-abelian simple group $T$ with an outer involutory automorphism $\alpha$ and $x\in T$ of odd order at least $5$ with $\langle x,\alpha(x)\rangle=T$, is $B(T;x,\alpha(x))$ always $3$-edge-colourable?  Same question for $T=S\times S$ with $\alpha$ the swap.  The order-three case is Proposition~\ref{prop:index2-three}; for larger odd orders, is the coset-incidence graph in Proposition~\ref{prop:consecutive} guaranteed to have a consecutive $2$-factor?  The case $T=A_n$, $\alpha$ conjugation by a transposition, is already interesting.
\item (Arc-regular cubic graphs.)  Is every connected cubic graph admitting an arc-regular group of automorphisms $3$-edge-colourable (Corollary~\ref{cor:arcregular})?  By Corollary~\ref{cor:conder} the answer is yes up to $10\,000$ vertices, and by Corollary~\ref{cor:three-simple} it is enough to exclude counterexamples whose arc-regular group is non-abelian simple.
\item (Vertex-transitive graphs.)  Are the Petersen graph and its truncation the only connected cubic vertex-transitive graphs that are not $3$-edge-colourable?  Corollary~\ref{cor:census} says yes up to $1280$ vertices, and Theorem~\ref{thm:potocnik} says yes for solvable vertex-transitive groups.
\item (Girth.)  Does Conjecture~\ref{conj:AZ} hold for cubic Cayley graphs of girth $3$ and $5$ (equivalently, by Proposition~\ref{prop:girth}, for generators $x$ of order $3$ and $5$), and of girth $6$?  In the girth-three case Corollary~\ref{cor:three-simple} reduces the problem to non-abelian simple groups.  Does the conjecture hold for girth at least $7$, i.e.\ is there a cyclically $7$-edge-connected Cayley snark?
\item (Hamiltonicity.)  Is every $(2,s,3)$-Cayley graph with $|G|\equiv0\pmod 4$ and $s\equiv2\pmod4$ hamiltonian?  (This is the remaining case of Theorem~\ref{thm:2s3}; it does not affect Conjecture~\ref{conj:AZ}.)
\item (Structure.)  Is there an algebraic characterisation of the even $2$-factors of $\X{G}{a}{x}$, for instance in terms of a subgroup $H$ of odd order such that the colouring is $H$-invariant?  Our computations show that in all cases tried an $H$-invariant colouring exists for a subgroup $H$ of index at most $9\,240$ (Table~\ref{tab:computation}), often a point stabiliser or an involution centraliser; is there always an $H$-invariant colouring with $H$ of bounded index (say $|G:H|\le c\,|G|^{1/2}$), or with $H$ from a specific family such as the centraliser of an involution?
\item (Cycle-double-cover palette.)  For the generator-separated flow of Proposition~\ref{prop:separated-eight-flow}, can one choose the quotient flow, its lifts around the $x$-cycles, and one of the twelve transvection matchings so that the augmented linear system in Proposition~\ref{prop:linearised-matching} is consistent?  Can the duality argument in the cycle-double-cover proof be extended to these extra equations?  This is already open when $x$ has prime order at least $5$, where the required generator-separated flow always exists.  Propositions~\ref{prop:full-rank-palette} and~\ref{prop:parallel-matching} explain why full rank and the four parallel matchings do not by themselves advance the conjecture.
\item (Weighted cut code.)  Can the generator-separated flow be chosen so that the syndrome criterion \eqref{eq:syndrome-criterion} holds for one transvection matching?  Propositions~\ref{prop:cdc-rigidity} and~\ref{prop:transvection-syndrome} show that the map from $\mathcal K_f$ to the dangerous-edge equations is exactly the effective freedom available after the flow and local triangles have been fixed.
\end{enumerate}

\subsection*{Acknowledgements}
This survey was prepared with the help of AI assistants (Claude and OpenAI Codex); the bibliographic details were checked against the cited sources where these were accessible, and the computations of Section~\ref{sec:computation} can be reproduced with the accompanying code.

\begin{thebibliography}{99}

\bibitem{AboomahigirNedela19} E.~Aboomahigir, R.~Nedela, Cubic Cayley graphs of girth at most six and their hamiltonicity, \emph{Acta Math. Univ. Comenian.} 88 (2019) 351--359.

\bibitem{AhanjidehIranmanesh19} M.~Ahanjideh, A.~Iranmanesh, The validity of Tutte's 3-flow conjecture for some Cayley graphs, \emph{Ars Math. Contemp.} 16 (2019) 203--213.

\bibitem{AhanjidehKovacs26} M.~Ahanjideh, I.~Kov\'acs, Nowhere-zero 3-flows in Cayley graphs on solvable groups of twice square-free order, arXiv:2603.24175 (2026).

\bibitem{ALZ96} B.~Alspach, Y.-P.~Liu, C.-Q.~Zhang, Nowhere-zero 4-flows and Cayley graphs on solvable groups, \emph{SIAM J. Discrete Math.} 9 (1996) 151--154.

\bibitem{CastagnaPrins72} F.~Castagna, G.~Prins, Every generalized Petersen graph has a Tait coloring, \emph{Pacific J. Math.} 40 (1972) 53--58.

\bibitem{ComptonWilliamson93} R.~C.~Compton, S.~G.~Williamson, Doubly adjacent Gray codes for the symmetric group, \emph{Linear Multilinear Algebra} 35 (1993) 237--293.

\bibitem{CaDiCaL} A.~Biere, M.~Fleury, N.~Froleyks, M.~J.~H.~Heule, CaDiCaL 2.0, in: \emph{Computer Aided Verification (CAV 2024)}, Lecture Notes in Comput. Sci. 14681, Springer, 2024, pp.~133--152.

\bibitem{ATLAS} R.~A.~Wilson et al., Atlas of Finite Group Representations, version 3, \url{https://brauer.maths.qmul.ac.uk/Atlas/v3/}.

\bibitem{Conder11} M.~Conder, Trivalent (cubic) symmetric graphs on up to 10000 vertices, 2011, \url{https://www.math.auckland.ac.nz/~conder/symmcubic10000list.txt}.

\bibitem{ConderNedela09} M.~Conder, R.~Nedela, A refined classification of symmetric cubic graphs, \emph{J. Algebra} 322 (2009) 722--740.

\bibitem{CurranGallian96} S.~J.~Curran, J.~A.~Gallian, Hamiltonian cycles and paths in Cayley graphs and digraphs---a survey, \emph{Discrete Math.} 156 (1996) 1--18.

\bibitem{DjokovicMiller80} D.~\v{Z}.~Djokovi\'c, G.~L.~Miller, Regular groups of automorphisms of cubic graphs, \emph{J. Combin. Theory Ser. B} 29 (1980) 195--230.

\bibitem{GKM09} H.~Glover, K.~Kutnar, D.~Maru\v{s}i\v{c}, Hamiltonian cycles in cubic Cayley graphs: the $\langle 2,4k,3\rangle$ case, \emph{J. Algebraic Combin.} 30 (2009) 447--475.

\bibitem{GKMM12} H.~Glover, K.~Kutnar, A.~Malni\v{c}, D.~Maru\v{s}i\v{c}, Hamilton cycles in (2,odd,3)-Cayley graphs, \emph{Proc. London Math. Soc.} (3) 104 (2012) 1171--1197.

\bibitem{GM07} H.~Glover, D.~Maru\v{s}i\v{c}, Hamiltonicity of cubic Cayley graphs, \emph{J. Eur. Math. Soc.} 9 (2007) 775--787.

\bibitem{Geelen26} J.~Geelen, OpenAI's proof of the Cycle Double Cover Theorem, arXiv:2607.15399 (2026).

\bibitem{HKM13} A.~Hujdurovi\'c, K.~Kutnar, D.~Maru\v{s}i\v{c}, On prime-valent symmetric bicirculants and Cayley snarks, in: \emph{Geometric Science of Information}, Lecture Notes in Comput. Sci. 8085, Springer, 2013, pp.~196--203.

\bibitem{HKM14} A.~Hujdurovi\'c, K.~Kutnar, D.~Maru\v{s}i\v{c}, Cubic Cayley graphs and snarks, in: R.~Connelly, A.~Ivi\'c Weiss, W.~Whiteley (eds.), \emph{Rigidity and Symmetry}, Fields Inst. Commun. 70, Springer, New York, 2014, pp.~27--40.

\bibitem{Jaeger79} F.~Jaeger, Flows and generalized coloring theorems in graphs, \emph{J. Combin. Theory Ser. B} 26 (1979) 205--216.

\bibitem{Jaeger85} F.~Jaeger, A survey of the cycle double cover conjecture, in: \emph{Cycles in Graphs}, Ann. Discrete Math. 27 (1985) 1--12.

\bibitem{JaegerSwart80} F.~Jaeger, T.~Swart, Problem session, in: \emph{Combinatorics 79}, Ann. Discrete Math. 9 (1980) 305.

\bibitem{Kochol96} M.~Kochol, Snarks without small cycles, \emph{J. Combin. Theory Ser. B} 67 (1996) 34--47.

\bibitem{KratochvilPoljak92} J.~Kratochv\'il, S.~Poljak, Compatible $2$-factors, \emph{Discrete Appl. Math.} 36 (1992) 253--266.

\bibitem{KMMMS12} K.~Kutnar, D.~Maru\v{s}i\v{c}, D.~W.~Morris, J.~Morris, P.~\v{S}parl, Hamiltonian cycles in Cayley graphs whose order has few prime factors, \emph{Ars Math. Contemp.} 5 (2012) 27--71.

\bibitem{MorrisWilk20} D.~W.~Morris, K.~Wilk, Cayley graphs of order $kp$ are hamiltonian for $k<48$, \emph{Art Discrete Appl. Math.} 3 (2020) \#P2.02; arXiv:1805.00149.

\bibitem{KutnarMarusic09} K.~Kutnar, D.~Maru\v{s}i\v{c}, Hamilton cycles and paths in vertex-transitive graphs---current directions, \emph{Discrete Math.} 309 (2009) 5491--5500.

\bibitem{LiZhou16} X.~Li, S.~Zhou, Nowhere-zero 3-flows in graphs admitting solvable arc-transitive groups of automorphisms, \emph{Ars Math. Contemp.} 10 (2016) 85--90.

\bibitem{LiebeckShalev96} M.~W.~Liebeck, A.~Shalev, Classical groups, probabilistic methods and the $(2,3)$-generation problem, \emph{Ann. of Math.} 144 (1996) 77--125.

\bibitem{Mader71} W.~Mader, Minimale $n$-fach kantenzusammenh\"angende Graphen, \emph{Math. Ann.} 191 (1971) 21--28.

\bibitem{MMP04} A.~Malni\v{c}, D.~Maru\v{s}i\v{c}, P.~Poto\v{c}nik, On cubic graphs admitting an edge-transitive solvable group, \emph{J. Algebraic Combin.} 20 (2004) 99--113.

\bibitem{MalnicPotocnik06} A.~Malni\v{c}, P.~Poto\v{c}nik, Invariant subspaces, duality, and covers of the Petersen graph, \emph{European J. Combin.} 27 (2006) 971--989.

\bibitem{NanasiovaSkoviera09} M.~N\'an\'asiov\'a, M.~\v{S}koviera, Nowhere-zero flows in Cayley graphs and Sylow 2-subgroups, \emph{J. Algebraic Combin.} 30 (2009) 103--110.

\bibitem{NS95} R.~Nedela, M.~\v{S}koviera, Atoms of cyclic connectivity in cubic graphs, \emph{Math. Slovaca} 45 (1995) 481--499.

\bibitem{NS95b} R.~Nedela, M.~\v{S}koviera, Which generalized Petersen graphs are Cayley graphs?, \emph{J. Graph Theory} 19 (1995) 1--11.

\bibitem{NS01} R.~Nedela, M.~\v{S}koviera, Cayley snarks and almost simple groups, \emph{Combinatorica} 21 (2001) 583--590.

\bibitem{OpenAICDC26} OpenAI, A proof of the cycle double cover conjecture, 2026, \url{https://cdn.openai.com/pdf/04d1d1e4-bc75-476a-97cf-49055cd98d31/cdc_proof.pdf}.

\bibitem{Oum26} S.-i.~Oum, A proof of the cycle double cover conjecture by OpenAI: an exposition, arXiv:2607.16356 (2026).

\bibitem{PakRadoicic09} I.~Pak, R.~Radoi\v{c}i\'c, Hamiltonian paths in Cayley graphs, \emph{Discrete Math.} 309 (2009) 5501--5508.

\bibitem{Pellegrini15} M.~A.~Pellegrini, M.~C.~Tamburini Bellani, The simple classical groups of dimension less than 6 which are $(2,3)$-generated, \emph{J. Algebra Appl.} 14 (2015) 1550148; arXiv:1405.3149.

\bibitem{Potocnik04} P.~Poto\v{c}nik, Edge-colourings of cubic graphs admitting a solvable vertex-transitive group of automorphisms, \emph{J. Combin. Theory Ser. B} 91 (2004) 289--300.

\bibitem{PSS05} P.~Poto\v{c}nik, M.~\v{S}koviera, R.~\v{S}krekovski, Nowhere-zero 3-flows in abelian Cayley graphs, \emph{Discrete Math.} 297 (2005) 119--127.

\bibitem{PSV13} P.~Poto\v{c}nik, P.~Spiga, G.~Verret, Cubic vertex-transitive graphs on up to 1280 vertices, \emph{J. Symbolic Comput.} 50 (2013) 465--477.

\bibitem{SawadaWilliams20} J.~Sawada, A.~Williams, Solving the sigma-tau problem, \emph{ACM Trans. Algorithms} 16 (2020) Art.~11.

\bibitem{Tutte54} W.~T.~Tutte, A contribution to the theory of chromatic polynomials, \emph{Canad. J. Math.} 6 (1954) 80--91.

\bibitem{Tutte66} W.~T.~Tutte, On the algebraic theory of graph colorings, \emph{J. Combin. Theory} 1 (1966) 15--50.

\bibitem{Woldar89} A.~J.~Woldar, On Hurwitz generation and genus actions of sporadic groups, \emph{Illinois J. Math.} 33 (1989) 416--437.

\bibitem{YangLi11} F.~Yang, X.~Li, Nowhere-zero 3-flows in dihedral Cayley graphs, \emph{Inform. Process. Lett.} 111 (2011) 416--419.

\bibitem{ZhangZhou22b} J.~Zhang, S.~Zhou, Nowhere-zero 3-flows in nilpotently vertex-transitive graphs, arXiv:2203.13575 (2022).

\bibitem{ZhangZhou24} J.~Zhang, S.~Zhou, Nowhere-zero 3-flows in Cayley graphs on supersolvable groups, \emph{J. Combin. Theory Ser. A} 204 (2024) 105852; arXiv:2203.02971.

\end{thebibliography}

\end{document}
